% Options for packages loaded elsewhere
\PassOptionsToPackage{unicode}{hyperref}
\PassOptionsToPackage{hyphens}{url}
\documentclass[
]{article}
\usepackage{xcolor}
\usepackage{amsmath,amssymb}
\setcounter{secnumdepth}{-\maxdimen} % remove section numbering
\usepackage{iftex}
\ifPDFTeX
  \usepackage[T1]{fontenc}
  \usepackage[utf8]{inputenc}
  \usepackage{textcomp} % provide euro and other symbols
\else % if luatex or xetex
  \usepackage{unicode-math} % this also loads fontspec
  \defaultfontfeatures{Scale=MatchLowercase}
  \defaultfontfeatures[\rmfamily]{Ligatures=TeX,Scale=1}
\fi
\usepackage{lmodern}
\ifPDFTeX\else
  % xetex/luatex font selection
\fi
% Use upquote if available, for straight quotes in verbatim environments
\IfFileExists{upquote.sty}{\usepackage{upquote}}{}
\IfFileExists{microtype.sty}{% use microtype if available
  \usepackage[]{microtype}
  \UseMicrotypeSet[protrusion]{basicmath} % disable protrusion for tt fonts
}{}
\makeatletter
\@ifundefined{KOMAClassName}{% if non-KOMA class
  \IfFileExists{parskip.sty}{%
    \usepackage{parskip}
  }{% else
    \setlength{\parindent}{0pt}
    \setlength{\parskip}{6pt plus 2pt minus 1pt}}
}{% if KOMA class
  \KOMAoptions{parskip=half}}
\makeatother
\usepackage{color}
\usepackage{fancyvrb}
\newcommand{\VerbBar}{|}
\newcommand{\VERB}{\Verb[commandchars=\\\{\}]}
\DefineVerbatimEnvironment{Highlighting}{Verbatim}{commandchars=\\\{\}}
% Add ',fontsize=\small' for more characters per line
\newenvironment{Shaded}{}{}
\newcommand{\AlertTok}[1]{\textcolor[rgb]{1.00,0.00,0.00}{\textbf{#1}}}
\newcommand{\AnnotationTok}[1]{\textcolor[rgb]{0.38,0.63,0.69}{\textbf{\textit{#1}}}}
\newcommand{\AttributeTok}[1]{\textcolor[rgb]{0.49,0.56,0.16}{#1}}
\newcommand{\BaseNTok}[1]{\textcolor[rgb]{0.25,0.63,0.44}{#1}}
\newcommand{\BuiltInTok}[1]{\textcolor[rgb]{0.00,0.50,0.00}{#1}}
\newcommand{\CharTok}[1]{\textcolor[rgb]{0.25,0.44,0.63}{#1}}
\newcommand{\CommentTok}[1]{\textcolor[rgb]{0.38,0.63,0.69}{\textit{#1}}}
\newcommand{\CommentVarTok}[1]{\textcolor[rgb]{0.38,0.63,0.69}{\textbf{\textit{#1}}}}
\newcommand{\ConstantTok}[1]{\textcolor[rgb]{0.53,0.00,0.00}{#1}}
\newcommand{\ControlFlowTok}[1]{\textcolor[rgb]{0.00,0.44,0.13}{\textbf{#1}}}
\newcommand{\DataTypeTok}[1]{\textcolor[rgb]{0.56,0.13,0.00}{#1}}
\newcommand{\DecValTok}[1]{\textcolor[rgb]{0.25,0.63,0.44}{#1}}
\newcommand{\DocumentationTok}[1]{\textcolor[rgb]{0.73,0.13,0.13}{\textit{#1}}}
\newcommand{\ErrorTok}[1]{\textcolor[rgb]{1.00,0.00,0.00}{\textbf{#1}}}
\newcommand{\ExtensionTok}[1]{#1}
\newcommand{\FloatTok}[1]{\textcolor[rgb]{0.25,0.63,0.44}{#1}}
\newcommand{\FunctionTok}[1]{\textcolor[rgb]{0.02,0.16,0.49}{#1}}
\newcommand{\ImportTok}[1]{\textcolor[rgb]{0.00,0.50,0.00}{\textbf{#1}}}
\newcommand{\InformationTok}[1]{\textcolor[rgb]{0.38,0.63,0.69}{\textbf{\textit{#1}}}}
\newcommand{\KeywordTok}[1]{\textcolor[rgb]{0.00,0.44,0.13}{\textbf{#1}}}
\newcommand{\NormalTok}[1]{#1}
\newcommand{\OperatorTok}[1]{\textcolor[rgb]{0.40,0.40,0.40}{#1}}
\newcommand{\OtherTok}[1]{\textcolor[rgb]{0.00,0.44,0.13}{#1}}
\newcommand{\PreprocessorTok}[1]{\textcolor[rgb]{0.74,0.48,0.00}{#1}}
\newcommand{\RegionMarkerTok}[1]{#1}
\newcommand{\SpecialCharTok}[1]{\textcolor[rgb]{0.25,0.44,0.63}{#1}}
\newcommand{\SpecialStringTok}[1]{\textcolor[rgb]{0.73,0.40,0.53}{#1}}
\newcommand{\StringTok}[1]{\textcolor[rgb]{0.25,0.44,0.63}{#1}}
\newcommand{\VariableTok}[1]{\textcolor[rgb]{0.10,0.09,0.49}{#1}}
\newcommand{\VerbatimStringTok}[1]{\textcolor[rgb]{0.25,0.44,0.63}{#1}}
\newcommand{\WarningTok}[1]{\textcolor[rgb]{0.38,0.63,0.69}{\textbf{\textit{#1}}}}
\usepackage{longtable,booktabs,array}
\newcounter{none} % for unnumbered tables
\usepackage{calc} % for calculating minipage widths
% Correct order of tables after \paragraph or \subparagraph
\usepackage{etoolbox}
\makeatletter
\patchcmd\longtable{\par}{\if@noskipsec\mbox{}\fi\par}{}{}
\makeatother
% Allow footnotes in longtable head/foot
\IfFileExists{footnotehyper.sty}{\usepackage{footnotehyper}}{\usepackage{footnote}}
\makesavenoteenv{longtable}
\setlength{\emergencystretch}{3em} % prevent overfull lines
\providecommand{\tightlist}{%
  \setlength{\itemsep}{0pt}\setlength{\parskip}{0pt}}
\usepackage{hyperxmp}
\usepackage{bookmark}
\IfFileExists{xurl.sty}{\usepackage{xurl}}{} % add URL line breaks if available
\urlstyle{same}
\hypersetup{
  pdftitle={Reactions and the partition: opt-in eventual consistency in actor-native systems},
  pdfauthor={Alvaro Rivera},
  pdfkeywords={\xmpquote{reactions}, \xmpquote{followers}, \xmpquote{consistency
contract}, \xmpquote{opt-in eventual consistency}, \xmpquote{domain
purity}, \xmpquote{closability}, \xmpquote{actor
model}, \xmpquote{CQRS}, \xmpquote{event
sourcing}, \xmpquote{placement}, \xmpquote{choreography}, \xmpquote{puppeteer
framework}},
  hidelinks,
  pdfcreator={LaTeX via pandoc}}

\title{Reactions and the partition: opt-in eventual consistency in
actor-native systems}
\author{Alvaro Rivera}
\date{2026-05-12}

\begin{document}
\maketitle
\begin{abstract}
This paper makes a design theory contribution in the sense of Hevner,
March, Park, and Ram (2004): it identifies a new design construct ---
the developer-controlled pragmatic partition of an event's implied work
into now and deferred --- derives the design principles required for
that partition to be exercised in practice, and presents an
instantiation in which those principles have been realized in
production. Eventual consistency is treated, in the CQRS literature, as
a defining commitment of the architecture: separated read and write
stores synchronize asynchronously, and applications must absorb the
contract. This paper argues for actor-native systems that the contract
is upside-down. The primitive is not eventual consistency; it is the
developer's pragmatic partition of the work an event implies into a
\emph{now} branch (what the verb's contract promises before responding)
and a \emph{deferred} branch (placed wherever the developer chose ---
same thread, another node, the cluster). Reactions are the artifact
through which that partition is exercised, and they are not async events
under another name: each Reaction matches a \emph{trajectory} through
the actor's journal, statically compiled against the domain's libraries,
with typed identifiers propagated across stages. Reactions serve a
double purpose --- pragmatic deferral and categorical segregation of
operational concerns away from domain methods --- and the segregation is
structural, not stylistic: the language itself blocks the operational
concern from leaking back into the actor's state. Both purposes share a
single guardian --- friction --- and collapse together if it is high.
Three consequences fall out of the primitive when its exercise is
low-friction: a domain library that closes (the primary consequence ---
declarable as ``task done'' because it is free of present and future
operational tools), fast verbs by construction, and opt-in eventual
consistency (the developer names each deferred effect by hand, inverting
the CQRS-classical assumption). The argument is conceptual; code
references throughout point to the public Puppeteer codebase, in which
the Reaction primitive --- its two modes (\texttt{Cue} and
\texttt{Job}), its three action planes (\texttt{Program.Emit},
\texttt{Causation.Continue}, \texttt{Metadata.Elide}), its activation
scopes, and its multi-event pattern matching --- instantiates the
primitive concretely.
\end{abstract}

\section{Reactions and the partition}\label{reactions-and-the-partition}

\subsection{TL;DR}\label{tldr}

\begin{quote}
When CQRS introduces eventual consistency, the standard treatment is as
a cost: a contract between separated read and write stores that the
application must absorb. This paper argues that for actor-native systems
the contract is upside-down. The primitive is not eventual consistency
at all; it is the developer's pragmatic partition of the work an event
implies into \emph{now} (what the verb's contract promises before
responding) and \emph{deferred} (placed wherever the developer chose ---
same thread, another node, the cluster).

Reactions are the artifact through which the partition is exercised ---
and they are not async events under another name. Each Reaction matches
a \emph{trajectory} through the actor's journal, a sequence of one or
more stages naming a saga or lifecycle, statically compiled against the
domain's libraries with typed identifiers propagated across stages.
Reactions serve two purposes simultaneously. The first is
\textbf{pragmatic deferral} --- for latency budget, parallelization,
scaling. The second is \textbf{categorical segregation} --- keeping
operational concerns (email, BI, third-party messaging, telemetry) from
propagating into domain methods. Reactions are \emph{not}
technology-free; they are precisely where that technology lives,
\textbf{encapsulated} so that other domain verbs do not inherit the
dependency, and the encapsulation is structural --- the language itself
blocks the operational concern from leaking back. The library, in turn,
stays free of that propagation, built under its legitimate programming
paradigm (currently object orientation). Both purposes share one
guardian: \textbf{friction}. If Reactions cost the developer more than
inline code, the partition is not exercised; verbs degrade and the
domain library is contaminated by precipitation. The friction is an
architectural responsibility of the framework, not a discipline question
for the developer.

Three consequences fall out of the primitive when it is low-friction.
The primary one is a \textbf{domain library that closes} --- that can be
declared \emph{task done} because it is structurally free of present and
future operational tools. The other two are \textbf{fast verbs by
construction} and \textbf{opt-in eventual consistency} (the developer
names each deferred effect by hand, inverting the CQRS-classical
assumption). \textbf{Eventual consistency is the shadow of the
partition, not the design choice.}
\end{quote}

\begin{center}\rule{0.5\linewidth}{0.5pt}\end{center}

\subsection{Claims this paper makes}\label{claims-this-paper-makes}

\begin{enumerate}
\def\labelenumi{\arabic{enumi}.}
\item
  \textbf{Logical, not physical, CQRS.} Within an actor, write and read
  paths operate on a single in-memory state. \texttt{PerformCmd} and
  \texttt{PerformQry} are verbs over the same store; there is no
  synchronization between separated write and read sides because there
  are no separated sides. The eventual-consistency contract that defines
  classical CQRS does not apply at the actor boundary.
  \emph{(Verification: code references to
  \texttt{PerformCmd}/\texttt{PerformQry} over a shared actor state in
  §6.)}
\item
  \textbf{Within an event's frame, formal correctness places everything
  in \emph{now}.} The work an event implies --- every domain effect
  entailed by the verb's contract --- should, formally, resolve before
  the event closes. The \emph{now/deferred} split is therefore not a
  decomposition deduced from correctness; it is a pragmatic boundary the
  developer draws inside the event's implied work. (``Everything'' here
  scopes to the event's domain-implied work; concerns extrinsic to the
  domain are addressed in claim 4.)
\item
  \textbf{The now/deferred partition is the primitive.} The developer
  chooses, feature by feature, what the verb's response promises now and
  what is deferred. The choice is pragmatic: it is governed by latency
  budget, parallelization opportunity, and the categorical purity
  argument made in claim 4 --- not by a deduction from the formal
  contract.
\item
  \textbf{The artifact realizing the partition serves a double purpose.}
  In the instantiation developed here, this artifact is the Reaction.

  \begin{itemize}
  \tightlist
  \item
    \textbf{(a) Pragmatic partition} --- work that the verb cannot
    afford to consume within its latency budget, or that is more
    naturally executed in parallel, is deferred.
  \item
    \textbf{(b) Categorical segregation} --- extrinsic operational
    concerns (email, BI, third-party messaging, telemetry, webhooks)
    demanded by requirements but not part of the rich library land in
    the deferral artifact, not in domain methods. The segregation is of
    \emph{propagation}, not of presence: the artifact may itself invoke
    email or telemetry technology, and that is appropriate --- the
    technology is encapsulated where it can be exercised without leaking
    into other domain verbs. \emph{The artifact is not technology-free;
    it is the encapsulation that keeps technology from propagating into
    the domain.} This extends the anti-porosity principle of
    \href{01-anti-porosity.md}{Paper 1} to a fourth layer --- the
    operational edge of the domain --- and is compatible with the rich
    library being built under the legitimate programming paradigm
    (currently object orientation) that sustains the domain language.
    That paradigm and the conceptual primitives it permits (classes,
    inheritance, polymorphism, invariants, domain methods) are the
    legitimate tool \emph{inside} the library (§3.1). \emph{The
    encapsulation is structurally enforced, not conventional: in
    Puppeteer's instantiation, a Reaction's emit path runs as a query
    (\texttt{PerformEmit} with \texttt{isQuery:true}, blocking
    \texttt{expose} and journal writes) --- the language refuses to let
    the operational concern leak back into the actor's state.
    (Verification: \texttt{ActorHandler.cs:1662-1733}.)}
  \end{itemize}
\item
  \textbf{Deferred work carries a placement decision.} The developer, at
  the moment of deferral, also decides where the deferred work runs
  along a \emph{proximity / urgency} axis: close-to-action (same thread,
  near-continuation, parent actor's memory available) ↔
  remote-and-deferred (another thread, another node,
  cluster-distributed). Placement is part of the primitive, not an
  implementation detail; it is orthogonal to whether the deferral is
  pragmatic (4a) or categorical (4b).
\item
  \textbf{Friction is the structural guardian of the partition.} The
  artifact realizing the partition must minimize the cost of declaring
  deferred work, or both purposes collapse simultaneously: if deferring
  costs the developer more lines or more cognitive load than inline
  code, the partition is not exercised, verbs degrade \emph{and} the
  domain library is contaminated by precipitation. The contamination is
  the worse of the two failures because it is hard to revert and
  propagates couplings. Friction is an architectural responsibility of
  any complete instantiation, not a question of developer discipline.
  \emph{(Verbatim author statement anchored as blockquote in §4.)}
\item
  \textbf{Three consequences of (3) + (4) + (6), ranked.} The primary
  consequence is \textbf{a domain library that closes} --- declarable as
  \emph{task done} by being free of present and future operational tools
  (§5.1, §10). Secondary consequences are \textbf{fast verbs by
  construction} (§5.2) and \textbf{opt-in eventual consistency}, in
  which the developer names each deferred effect by hand --- the inverse
  of the CQRS-classical opt-out assumption (§5.3, honors the
  forward-reference of \href{01-anti-porosity.md}{Paper 1 §6}). The
  three are not features and not design choices; they fall out of the
  primitive when it is low-friction. \emph{Fast verbs in particular are
  structural: a Reaction executes outside the write semaphore of
  \texttt{PerformCommand}; the verb returns to the caller without
  waiting for any deferred work to begin. (Verification:
  \texttt{Reaction.cs:715-832}, \texttt{ActorHandler.cs:1662-1733}.)}
\item
  \textbf{The continuation contract survives the deployment switch.}
  Deferred work in flight when a node hands off to its successor is not
  lost; red-black is orchestrated by the runtime --- developers do not
  normally see the mechanism. The only common exception is coordinated
  cut-over of a Kafka producer/consumer pair when a message format
  changes, where both sides must release together. In Puppeteer's
  instantiation, this property is realized by
  \texttt{Theater.aliveGate}; the mechanism itself is treated in detail
  in \href{05-zero-downtime.md}{Paper 5} (zero-downtime deployment,
  forthcoming). The cross-actor continuity primitive that complements
  the in-actor continuity argued here is the subject of
  \href{04-cross-actor-continuity.md}{Paper 4}, already published in the
  present series.
\item
  \textbf{Domain libraries become real artifacts.} In most enterprises
  today, ``the domain'' does not exist as a library --- it lives as
  classes scattered across each csproj, copied between services,
  drifting independently over time until ``the customer model'' or ``the
  order model'' has become a dozen incompatible structures across the
  codebase. The pattern of pure domain methods, preserved by Reactions
  encapsulating operational concerns (claim 4b), is the structural
  condition under which a single domain library, maintained as one
  artifact over years, becomes economically possible at all.
  \emph{Reusability across endpoints is the operational form of
  closability} (claim 7).
\item
  \textbf{The artifact realizing the partition is declarative,
  multi-event pattern matching against the actor's trajectory,
  statically compiled against the domain's libraries.} It is not an
  event handler. It is a rule the engine applies over the actor's
  journal: a sequence of one or more stages that, taken together,
  describe a trajectory through the actor's history --- a saga, a
  lifecycle, a recurring anomaly. Identifiers captured in one stage
  propagate to the next, correlating the events. The pattern is parsed
  and resolved against the domain's libraries in build phase, not
  against text at runtime; a stage referring to a class the domain does
  not contain fails to compile. The propagation, the typed capture, and
  the static binding to the domain are jointly what distinguishes this
  primitive from event-driven continuations in Akka, grain reminders in
  Orleans, or topic consumers in Kafka. The instantiation that
  demonstrates this is the Reaction primitive of Puppeteer --- its
  \texttt{Seek}, \texttt{ThenSeek}, and \texttt{ThenFinalSeek} stages
  are developed in §6.2 and §6.3 \emph{(primary code anchors
  \texttt{Pattern.cs}, \texttt{Reaction.SolveActionReferences()} at
  \texttt{Reaction.cs:997}, \texttt{MatchTree.TryMatchAtLevel()})}.
\item
  \textbf{The proximity/urgency axis has two canonical points, not a
  continuous spectrum.} In Puppeteer's instantiation, these are
  \texttt{Cue()} (push loop on a separate thread, operating against the
  actor's live in-memory state, \texttt{ReadForward()} mandatory,
  sub-second latency) and \texttt{Job()} (on-demand against the
  replicated journal, no co-location with the live actor, freely
  placeable on a follower in another pod or another machine in the
  cluster). The two modes are structurally distinct, not points on a
  slider --- they differ in source of state (memory vs journal), in
  trigger (push callback vs explicit invocation), and in what they
  presuppose about the host. The placement decision of claim 5 is not
  abstract: it is the choice of one mode versus the other, possibly
  further narrowed by activation scope (in Puppeteer:
  \texttt{DirectorOnly}, \texttt{CastOnly}, \texttt{Company}).
  \emph{(Verification: §6.1 and §6.6; primary code anchor
  \texttt{Reaction.cs:19-23} and \texttt{Reaction.cs:67-80}.)}
\item
  \textbf{The artifact's elision action integrates declarative
  trajectory correlation with journal density.} When the pattern
  correlates a multi-event trajectory (an order's
  \texttt{Initiate\ →\ Confirm\ →\ Pay} sequence, for instance), the
  runtime can mark the constituent events as elided in the journal.
  Subsequent rehydrations skip past them: the trajectory is closed, its
  component events have no further bearing on any decision. This closes
  a conceptual loop with \href{01-anti-porosity.md}{Paper 1} --- the
  \emph{skips} introduced there as a journal-density mechanism become,
  in the present construct, the natural outcome of a declarative
  correlation rather than a post-hoc optimization. In Puppeteer this is
  realized as \texttt{Metadata.Elide()} on the Reaction's metadata plane
  (§6.5; primary code anchor \texttt{Reaction.cs:715-832} and
  \texttt{DiaryStorage.EventElisionStorage.MarkEventsAsElided}).
\end{enumerate}

\subsection{When NOT to use this
approach}\label{when-not-to-use-this-approach}

The primitive of pragmatic partition is general; the implementation
pattern that realizes it has narrower limits. We name three regimes in
which the analysis presented here either does not apply, applies only
partially, or invites overreach.

\subsubsection{A. Domains where every effect must be transactionally
atomic with the
response}\label{a.-domains-where-every-effect-must-be-transactionally-atomic-with-the-response}

Some domains --- banking core ledgers, certain regulated transaction
systems --- require that every observable effect of a verb commit in
lockstep with the response. The pragmatic partition cannot be exercised
here without violating the domain's own contract. The opt-in framing
offered in claim 7 does not soften this: opt-in is a license the
developer extends to a feature; if the feature has no license to extend,
the option is null.

\subsubsection{B. Teams without appetite to think in terms of partition
decisions}\label{b.-teams-without-appetite-to-think-in-terms-of-partition-decisions}

The primitive shifts a design responsibility onto the developer that
some teams do not want. A team that treats Reactions as ``fire and
forget'' --- without modeling the contract that follows --- loses the
very property the primitive was supposed to confer. The framework can
lower the friction; it cannot supply the design judgment.

\subsubsection{C. Workflows where the deferred branch's failure is
structurally
inadmissible}\label{c.-workflows-where-the-deferred-branchs-failure-is-structurally-inadmissible}

A pragmatic partition presupposes that the deferred branch can fail or
retry without compromising the verb's already-committed contract.
Workflows in which a downstream failure of the deferred work must abort
the verb retroactively --- sagas with compensating actions are one such
pattern, but the right place to discuss them at depth is the companion
paper on Choreography. Some such workflows are admissible under the
partition; others are not. The honest position is that the partition is
not the right primitive for every long-running workflow.

\begin{center}\rule{0.5\linewidth}{0.5pt}\end{center}

\subsection{1. Introduction: a developer says
things}\label{introduction-a-developer-says-things}

This paper makes a design theory contribution. It identifies a
structural primitive that frameworks for actor-native systems have left
implicit --- the developer-controlled partition of an event's implied
work into a ``now'' branch and a ``deferred'' branch with a placement
decision attached --- and derives the design principles required to
exercise that primitive at scale. The instantiation that demonstrates
the construct is realizable is the Reaction primitive of the Puppeteer
framework, treated in §6. The genre is the one Hevner, March, Park, and
Ram (2004) name design science research: empirical evidence is presented
in the form of a working artifact rather than a controlled experiment.

A developer is modeling a purchase-order endpoint. They are not thinking
about transactions, consistency contracts, or CQRS sides. They are
thinking about roles and what each role does next: \emph{``the cashier
confirms the purchase and locks the requested items. Then a confirmation
goes out, and the rewarder evaluates whether any promotional campaign
applies.''} They say it at a whiteboard, and a few hours later it is in
code --- not as a set of architectural decisions translated through a
vocabulary, but as the same sentence the developer just spoke. In C\# a
developer says things --- defines classes, writes methods, encodes
invariants. What an actor-native runtime adds is that the same act of
saying things now extends to the endpoint level: who handles the
request, what the response promises, what is handled afterwards.

That extension carries a question the developer has not been asked to
answer in industrial practice for some time, and most of this paper is
about it. When the developer said \emph{``first this, then that, then
the other''}, they drew a partition --- between what the verb's response
promises \emph{now} and what is \emph{deferred and placed elsewhere}.
They drew it without naming it, on pragmatic grounds: latency,
parallelization, and the sense that some concerns simply do not belong
inside a domain method. The literature has names for the consequences of
that partition --- domain libraries that close, fast verbs by
construction, opt-in eventual consistency --- but the developer never
reached for those names. They reached for a sentence about who does
what, and the runtime was kind enough to let them write the sentence as
the program.

The classical actor model --- Akka, Orleans, Proto.Actor and their
lineage --- provides actor isolation and message passing, but it does
not provide a declarative, journal-respected mechanism for what the verb
commits to as part of its contract. What happens after the verb is, in
those frameworks, side effect: a method call into another grain, a
\texttt{tell} to another actor, a registered timer, a write to a queue.
None of those are part of the actor's program in the sense the journal
preserves (substrate-level \emph{program} throughout this paper refers
to the construct of Paper 2 §1.2: the pair of domain library and journal
of invocations). This paper argues that Reactions are that missing
mechanism --- and that without it, the journal-density property
established in \href{01-anti-porosity.md}{Paper 1} and
\href{02-program-value-separability.md}{Paper 2} cannot survive contact
with operational requirements that span the verb's response.

This paper takes that sentence apart in turn. §2 names the partition the
developer drew. §3 examines the artifacts the developer composed to
write it: a verb on the actor --- a command, a query, or a
check-then-command --- and the Reactions that handle what the verb
deferred. §4 looks at what makes writing a Reaction cheap or expensive,
and why the answer matters more than it appears at first. §5 follows
what falls out when the writing is cheap: three consequences, ranked. §6
grounds the discussion in code from the public Puppeteer codebase. §7
places the design against Akka and Orleans, §8 addresses the predictable
objections, §9 sets the work in the broader literature, and §10 returns
to the developer's sentence and observes what changed.

\subsection{2. The pragmatic partition: now versus
deferred}\label{the-pragmatic-partition-now-versus-deferred}

\subsubsection{\texorpdfstring{2.1 The first cut is
\emph{who}}{2.1 The first cut is who}}\label{the-first-cut-is-who}

The first cut the developer makes when modeling is not ``what feature is
this'' --- it is \emph{who}. Cashier. Packer. Deliver. Rewarder. Once
the actors are named, the work distributes itself: each role naturally
absorbs the verbs that belong to it, and the partition between
\emph{what I do now} and \emph{what I see go by and react to afterwards}
emerges from the modeling, not from a separate engineering decision. The
developer never says ``this is opt-in eventual consistency''; they say
what each role does. That sentence carries both role and feature --- the
cashier role \emph{and} the priority order, what is settled before
responding versus what is handled after --- and it is not
feature-by-feature; it is role-by-role-with-feature.

\subsubsection{2.2 The partition is drawn, not
deduced}\label{the-partition-is-drawn-not-deduced}

Formally, the situation is sharper. Every domain effect implied by the
event sits, by formal correctness alone, in \emph{now}: nothing in the
verb's contract entitles the developer to defer anything. The license to
defer is therefore not deduced --- it is \textbf{drawn}. The developer
draws it while modeling, on pragmatic grounds, and the grounds are
three. First, a latency budget the verb cannot exceed without breaking
the actor's serial-mailbox invariant --- the topic of §5.2. Second, an
opportunity to parallelize work across roles or across nodes --- the
placement question of §2.3. Third, a categorical sense that certain
concerns (a confirmation going out, an analytics ping, a third-party
update) do not belong inside a domain method at all, regardless of
timing --- the case developed in §3.3. The three grounds are independent
of one another, and a single deferral may rest on more than one.

\subsubsection{2.3 Placement is part of the same
act}\label{placement-is-part-of-the-same-act}

\emph{{[}Skeleton --- to iterate in chat. Refined 2026-05-06 to anchor
the axis on the two canonical points (Cue and Job) per claim 11.{]}} At
the same act of drawing the now/deferred partition, the developer also
decides where the deferred work runs. The axis is not continuous: it has
two canonical points named explicitly by the framework. \emph{Cue} runs
immediately on a separate thread, against the actor's live in-memory
state; \emph{Job} runs on demand against the replicated journal, freely
placeable on a follower in another pod or another machine. Section 6.1
develops the two modes in detail. The point relevant here is that
partition and placement are one decision, not two --- choosing what gets
deferred and choosing which of the two structurally distinct modes
carries it are the same act, not separable concerns.

This partition --- its drawing and its placement, taken as one act ---
is the primitive of the rest of the paper. Everything that follows is
downstream of it. §3 examines the artifacts through which the developer
writes the partition (a verb on the actor and the Reactions that handle
what the verb deferred), naming the two purposes those Reactions serve.
§4 looks at what makes writing them cheap or expensive. §5 follows what
falls out when the writing is cheap.

\subsection{3. The double purpose of
Reactions}\label{the-double-purpose-of-reactions}

\emph{{[}Skeleton.{]}} Before naming what Reactions exclude from the
domain library, the paper must name what they do \textbf{not} exclude.
The domain library is not tool-free. It is built under the legitimate
programming paradigm --- currently object orientation --- that sustains
the language of the domain itself. Once that legitimate tool is named
(§3.1), the two reasons to defer can be developed without inviting the
misreading that ``pure'' means ``tool-free'': pragmatic deferral (§3.2)
and categorical segregation of operational pollution (§3.3).

\subsubsection{3.1 The legitimate tool inside the
domain}\label{the-legitimate-tool-inside-the-domain}

The rich library is written in a host language under a programming
paradigm --- currently object orientation. Classes, inheritance,
polymorphism, invariants, domain methods, and the composition of
conceptual primitives are the \emph{legitimate tool} of the domain. They
live inside the library; they are what the library is made of. The
categorical segregation developed in §3.3 does not exclude this paradigm
--- on the contrary, it is what protects it.

The DSL of an actor runtime does not compete with this paradigm. It is
not a second implementation of the domain; it is the language by which
the domain is invoked. The relation is structurally analogous to that
between SQL and a relational engine: SQL does not duplicate the engine's
logic; it names operations the engine performs. A short SQL statement
can trigger joins, locks, and execution plans far heavier than its
surface text. The DSL of an actor runtime carries the same kind of
leverage --- a few lines invoke domain methods whose implementation
lives once in the library and whose behavior may be deep. Two endpoints
that share invocation lines are not duplicating knowledge; they are
composing distinct calls over the same library, in the sense Hunt and
Thomas (1999) named when they defined DRY as the prohibition on
duplicate \emph{representation of knowledge}, not on the appearance of
similar text. The polymorphism of the domain is a first-class citizen of
the DSL (\href{01-anti-porosity.md}{Paper 1} covers the parser and
symbol-table treatment of inheritance and interfaces).

A second observation belongs here, anticipating §5.1: the closability
argument that follows is about the boundary of the library, not about
its exhaustiveness. The library can keep growing in domain concepts
under its legitimate paradigm --- new classes, new invariants, new
relations --- and that growth is healthy and expected. Practical
incompleteness is not a defect of the closability argument. What closes
is the absorption of \emph{operational} tools that change with the
ecosystem; what remains open is the domain's own conceptual development.

\subsubsection{3.2 Pragmatic deferral}\label{pragmatic-deferral}

The first reason a developer defers work is operational. The verb of the
actor must complete promptly: an actor sustains its target throughput
when its mailbox advances quickly, and a verb that takes too long delays
every subsequent verb behind it. The constraint is not that I/O must be
absent from the verb --- formally permitted, an I/O may sit inside a
\texttt{PerformCommand} --- but that I/O accumulates badly. Many verbs
each making a network call, or one verb whose latency is unbounded, both
collapse the throughput the actor model is supposed to sustain. The fast
verb is structural rather than aspirational, and the developer protects
it by deferring.

Three patterns, each at a different scope, sustain the fast verb.
\emph{I/O hoisting at the caller} --- the caller resolves expensive
lookups and passes the resolved values into the \texttt{PerformCommand},
so the actor receives data and never fetches. \emph{Saga-phasing between
events} --- when a workflow genuinely requires interleaved external
calls, the actor advances one phase, releases, an external coordinator
does the I/O, and the response re-enters as the next event.
\emph{Reactions} --- work that the verb commits to as part of the
contract but does not block on, deferred to a \texttt{Cue} or
\texttt{Job} (§6.1). The three are not alternatives; they compose. What
§3.2 names is the third --- the case where the deferral is owned by the
actor itself, declared as a Reaction.

A practical recommendation belongs to the third pattern. A Reaction's
body is fast-verb code, with the same throughput constraint as any verb
on the actor: it is invoked exactly once per match. That guarantee ---
\emph{exactly once per match} --- is what permits a Reaction to declare
a side effect (sending a Kafka message, calling an external service)
inside what is otherwise pure domain code. Were the same side effect to
live inside a \texttt{PerformCommand}, it would violate the actor
principle that commands compute against in-memory state without external
I/O; lifting it into a Reaction keeps the actor's command path pure
while keeping the side effect's logic in the domain layer where it
belongs.

The same property --- \emph{exactly once per match} --- is what allows
the pattern to survive multi-datacenter replication, but only when the
side effect is mediated by a broker rather than performed directly
against shared infrastructure. Under journal replication, what travels
between datacenters is the action: its \texttt{actionId} and parameters.
Each datacenter runs the same domain code on its own actor, and each
datacenter's actor runs the same Reaction in turn, against its own
locally-replicated journal. If the Reaction emits to its \emph{local}
broker, and a \emph{local} consumer downstream writes to a \emph{local}
RDBMS or BI store, every datacenter behaves identically and
independently --- the topology is symmetric. If instead the Reaction
performs a synchronous write to a shared external RDBMS, both
datacenters now contend for the same external resource on every match,
and the operational independence the architecture was buying is lost.

The recommended pattern is therefore to keep slow I/O \emph{outside} the
Reaction body altogether: the Reaction emits a message to a broker ---
Kafka in the canonical case --- and a consumer downstream (a BI sink, an
analytics writer, a third-party adapter) performs whatever heavy I/O the
requirement demands. A typical BI consumer reduces to one insert and
several selects against its own store; the actor never sees that load.
The framework permits direct RDBMS calls inside a Reaction; the
recommendation is not to use them.

\subsubsection{3.3 Categorical segregation: encapsulation, not
exclusion}\label{categorical-segregation-encapsulation-not-exclusion}

Operational concerns demanded by requirements but extrinsic to the
domain --- email, BI, third-party messaging, telemetry, webhooks ---
land in Reactions, not in domain methods. The crucial point: the
segregation is of \emph{propagation}, not of \emph{presence}. A Reaction
can live inside the same library as the domain methods it observes ---
declared in the same project, next to the classes it cares about. What
the framework requires is that the Reaction be its own artifact, with
its own scope of effects, rather than collapsed into the body of a
\texttt{PerformCommand}. The structure is therefore one of co-location
without co-mingling. A Reaction whose purpose is to send a confirmation
email will, of course, contain a call to an email service; the
technology lives there, encapsulated, near the domain methods it serves
but distinct from them. What the segregation prevents is that the email
concern leak into a domain method that ought to know nothing about email
--- and that the next domain verb the actor processes inherit a
dependency on the email API by virtue of being in the same class.
\emph{Reactions are not technology-free; they are the encapsulation that
keeps technology from propagating into the domain.}

The function is structurally analogous to an anti-corruption layer in
DDD, but architectural rather than stylistic: any complete instantiation
treats the deferral artifact as a first-class citizen of the language,
and the developer's defaults route operational concerns there. In
Puppeteer this artifact is the Reaction. The legitimate paradigmatic
tool inside the library (§3.1) is preserved precisely because the
encapsulation is non-leaky.

The non-leak is structural, not stylistic. The Reaction's emit path runs
as a query: \texttt{PerformEmit} is invoked with \texttt{isQuery:true},
which blocks the script from using \texttt{expose} or declaring globals,
takes a read lock that runs in parallel with queries and other emits,
and leaves the journal untouched (§6.5; verification at
\texttt{ActorHandler.cs:1662-1733}). The technology of the Reaction (a
Kafka producer, an HTTP webhook, a BI sink) executes; the journal does
not record it as a domain event; the actor's state does not absorb the
call. The language refuses to let the operational concern leak back into
the actor in the parse phase, before the developer ever has the chance
to make a mistake about it.

Section 3.3 connects directly to \href{01-anti-porosity.md}{Paper 1}:
the anti-porosity transversal, treated there as domain + endpoint +
persistence, extends here to a fourth layer --- the operational edge of
the domain. \emph{Reactions are the anti-porous boundary at the
operational edge.}

\subsubsection{3.4 Orthogonality of purpose and
placement}\label{orthogonality-of-purpose-and-placement}

The two purposes named in §3.2 and §3.3 --- pragmatic deferral and
categorical segregation --- and the two placement modes named in §6.1
--- \texttt{Cue} and \texttt{Job} --- are independent decisions. A
pragmatic deferral may place close to the action as a \texttt{Cue}
near-continuation, or far from it as a \texttt{Job} running on a
journal-only follower; a categorical segregation may do the same. The
cells of the resulting two-by-two are all populated in practice: a
confirmation email may be a \texttt{Cue} (purpose categorical, placement
local) or a \texttt{Job} on a follower (purpose categorical, placement
remote); a heavy projection update may be a \texttt{Cue} against the
live actor (purpose pragmatic, placement local) or a \texttt{Job}
distributed across the cluster (purpose pragmatic, placement remote).
The developer chooses both dimensions at the moment of declaring the
artifact; the two dimensions remain independent in any complete
instantiation --- there is no ``default placement for categorical
concerns'' or ``obligatory mode for pragmatic deferrals.'' Each
declaration states its own combination.

A second observation about the partition belongs here. The partition is
not only a \emph{pragmatic} decision (latency, parallelization,
categorical separation); it is also a \emph{preservation} mechanism.
When the actor replicates across datacenters (§3.2), the journal travels
with it and so does the code: each datacenter executes the same
Reactions on its own actor against its own locally-replicated journal.
Were the verb's consequences side effects bolted onto the actor at
runtime --- \texttt{tell}s, registered timers, callbacks --- the
replication would either lose them or reproduce them incoherently.
Because Reactions are part of the program, declared in the DSL and
replicated as code, the actor behaves as its program says, in every
datacenter, every time it rehydrates.

\subsubsection{3.5 Pattern matching against the actor's
trajectory}\label{pattern-matching-against-the-actors-trajectory}

A Reaction is not an event handler. The framework's primitive is more
general --- and more difficult to find an analogue for in adjacent
runtimes. A Reaction is a rule the engine applies over the actor's
journal: a sequence of one or more \texttt{Seek} stages that, taken
together, describe a \emph{trajectory} through the actor's history. The
pattern matches not when one event arrives, but when the actor has, over
its lifetime, gone through the stages the rule names. Identifiers
captured in one stage propagate to the next, correlating the events: an
\texttt{order} named in the first \texttt{Seek} is the same
\texttt{order} that must appear in the next, and in the next. The
Reaction is part of the actor's program --- declared in the DSL, parsed
at build time against the domain's libraries, replicated as code
together with the journal --- not a side effect attached to the actor at
runtime.

Three properties are jointly characteristic. \emph{Multi-event}: the
smallest match is a saga, not a single message. \emph{Statically
compiled}: the pattern is parsed against the domain's \texttt{Libraries}
in build phase --- \texttt{{[}\_:Customer{]}.InitiateOrder(order:Order)}
resolves to the class \texttt{Customer} and to the signature of
\texttt{InitiateOrder}; if either does not exist in the domain, the
Reaction fails to construct. \emph{Trajectorial, not payload-based}:
what the matcher inspects is the sequence of \emph{invocations} the
actor recorded --- its conduct --- not the body of any one message.
These three together distinguish the primitive from event handlers in
Akka, grain reminders or streams in Orleans, and topic consumers in
Kafka. The closest analogue in the broader literature is complex event
processing, but CEP engines typically operate on external streams
without the host language's type system available; a Reaction operates
inside the actor's own type universe.

This is the property that licenses everything else. Without trajectorial
matching, \emph{``the cashier's verb is fast and a Reaction handles the
rewarder afterwards''} reduces to a callback. With it, what the
developer has written is \emph{``when, in the cashier's history, an
order was initiated, then confirmed, then paid, react''} --- a
declaration about how the actor's lifecycle unfolds, not about what to
do when a single event arrives. The matching machinery is examined in
§6.2; a worked example with the order lifecycle is in §6.3.

\subsubsection{3.6 Reactions are journal-indexed, not
event-driven}\label{reactions-are-journal-indexed-not-event-driven}

A Reaction is not an event handler. It is a journal-indexed program
cursor.

The structural property is this: \textbf{a Reaction traverses the
journal monotonically and exactly once.} No journal entry is ever
processed twice by the same Reaction --- not after a crash, not after a
restart, not after a replica handoff. Reprocessing is not something the
Reaction body avoids; it is something the construct forbids by shape.

The property does not rest on deduplication logic, idempotency wrappers,
or guards written inside the Reaction. It is upstream of all of those. A
Reaction's position over the journal --- one position per \texttt{Seek}
of the pattern --- is itself part of the actor's persistent state,
advances only forward, and is validated on load such that any
non-monotonic checkpoint is refused. The cursor is the Reaction's
identity over time; loss of the cursor is loss of the Reaction.

The runtime's commit discipline exists to sustain this property, not to
add it. The cursor's advance is durable before the Reaction's action
takes effect, and a process failure between those two moments leaves the
cursor ahead, never behind. The entry has already been past the cursor
by definition; on restart the Reaction's scan begins from a position
that excludes it. The asymmetry is deliberate: re-executing on the same
entry is the failure mode the construct is designed to forbid, while
failing to execute is the failure mode it is designed to expose ---
surfaced as a recoverability event rather than as silent re-execution.

The reader who is about to encounter cross-actor causation in
\href{04-cross-actor-continuity.md}{Paper 4} should hold this property
fixed. A \texttt{Causation.Continue} invocation that records a
cross-actor send in the sender's journal cannot fire twice against the
same journal entry --- not as a runtime accident, but as a consequence
of what a Reaction is. Tell turns the journal into the boundary of
causation precisely because the journal is a once-traversed program, not
a replayable log of handlers. The property is established here and
depended on there.

\subsection{4. Friction as the architectural
guardian}\label{friction-as-the-architectural-guardian}

The two purposes of §3 share a guardian. If Reactions cost the developer
more lines or more cognitive load than inline code, the partition is not
exercised. Both purposes collapse simultaneously: verbs degrade
\emph{and} the domain library is contaminated by precipitation.

\begin{quote}
\emph{Reactions must be at least as ergonomic as inline code. If
deferring costs the developer more than not deferring, the partition
isn't exercised --- verbs degrade AND the domain library is contaminated
by precipitation. The contamination is the worse failure: it is hard to
revert and propagates couplings. The friction is an architectural
responsibility of the framework, not a failure of the developer's
discipline.}
\end{quote}

The contamination is the worse of the two failures, for two reasons.
First, a slow verb is observable and gets fixed: monitoring catches it,
profiling diagnoses it, the team rewrites the slow verb on a known
schedule. A polluted domain method looks correct and does not --- its
tests pass, its code reads cleanly, the email it sends is the right
email. The dependency on the email API is invisible from the outside,
hidden behind the method's signature, until the day the developer needs
to revert it. Second, the contamination propagates. Once
\texttt{sendEmail(...)} lives inside a domain method, every caller of
that method inherits the coupling --- directly, by transitive call;
indirectly, by sharing a class with that method and inheriting its
dependencies in test setups, in mocking frameworks, in the surface area
of integration. Reverting the design later requires touching every call
site, every test, every place the class has appeared as an object of
composition.

The architectural responsibility is therefore not to teach the developer
discipline. It is to make Reactions, by construction, no more expensive
than inline code. The framework's contribution is not the partition
itself --- the developer always \emph{could} have deferred, in any actor
framework. The contribution is friction reduction along two dimensions.
The first is language-level cost: declaring a Reaction is a few lines of
fluent DSL, comparable to or shorter than an inline implementation of
the same effect (§6.1, §6.5). The second is textual proximity: Reactions
can be written next to the endpoint they observe --- in the same file, a
few lines below the \texttt{PerformCommand} they continue --- so that
the developer reading the endpoint sees, on the same screen, what else
happens around it: this confirmation goes out, that loyalty rule fires,
this analytics counter is updated. The Reaction's location is itself a
documentation aid. When the cost of doing the right thing exceeds the
cost of doing the wrong thing, no amount of design review will hold;
when the costs are equal or inverted --- and when the right thing is
also the more legible thing --- the partition exercises itself. The
framework that takes this seriously is the framework that gets the
partition for free.

\subsection{5. Three consequences
(ranked)}\label{three-consequences-ranked}

\emph{{[}Skeleton. The three consequences fall out of the primitive (§3)
when its exercise is low-friction (§4). They are ranked, with
closability primary, per the signed thesis of 2026-05-05. Subsections
5.1, 5.2, 5.3 in the order of importance argued by the rhetoric of the
series.{]}}

\subsubsection{5.1 A domain that closes}\label{a-domain-that-closes}

The primary consequence of the partition, when its exercise is
low-friction, is a domain library that closes. The library is not just
kept clean of operational pollution at any one moment; it is
structurally protected against the absorption of operational tools that
arrive in the future. Email today, Slack tomorrow, a webhook protocol
the year after --- none of these find their way into a domain method,
because the path of entry has been removed at the language level (§3.3).
The library that the team is writing now will still be the library the
team is reading in three years, and the conceptual primitives it
contains will still mean what they meant when they were declared.

\begin{quote}
\emph{At some point a domain must be declarable as ``task done''. A
library whose definition must absorb every present and future
operational tool --- email today, Slack tomorrow, a webhook protocol the
year after --- is a library that never closes. The only way it closes is
if it stays pure: free of operational concerns by construction.
Reactions are the artifact that makes that closure possible.}
\end{quote}

Two clarifications matter. First, what closes is the boundary, not the
conceptual growth of the domain itself. The library can keep absorbing
new domain concepts under its legitimate paradigm --- new classes, new
invariants, new relations --- and that growth is healthy and expected.
Practical incompleteness is not a defect of the closability argument:

\begin{quote}
\emph{If Puppeteer is not exhaustive, that is a practical limitation,
not a claim that the domain must be seen as complete.}
\end{quote}

Second, ``task done'' is an organisational sentence, not a
software-engineering one. A team that can declare its domain library
complete with respect to operational tools can build everything else ---
refactor, onboarding, upgrade, extension --- on a stable foundation. In
most enterprises today, ``the domain'' does not exist as a library at
all. It exists as classes scattered across each csproj, copied between
services, drifting independently over years until ``the customer model''
or ``the order model'' has become a dozen incompatible structures across
the codebase. The team speaks of these as if they were one thing; in
operational reality they are not. A library that closes its boundary to
operational tools while remaining open to conceptual growth is the
structural condition under which a single domain artifact, maintained as
one over years, becomes economically possible at all.

The closability argument extends the anti-porosity transversal of
\href{01-anti-porosity.md}{Paper 1} to a fourth layer. Domain libraries
shaped for the operations they describe and the conceptual primitives
they admit, not for the operational tools requirements happen to demand
this year, are libraries that the team can put down and pick back up
without forgetting what they meant.

\subsubsection{5.2 Fast verbs by
construction}\label{fast-verbs-by-construction}

Verbs in single-digit milliseconds are not a postulate of this
framework, but their explanation is layered. The actor's verbs are fast
first because the actor's state lives in memory: \texttt{PerformCmd} and
\texttt{PerformQry} operate on it without touching disk or external
stores. This is claim 1 of this paper --- logical CQRS --- and the
broader argument of \href{01-anti-porosity.md}{Paper 1}. They are fast
also because the domain methods invoked from a verb are themselves rich
but local: a single verb can dispatch many calls into the host language
and across many levels of the domain library, each operation cheap
because each operates on the same in-memory state
(\href{02-program-value-separability.md}{Paper 2}, in the empirical
measurements of the purchase example). The Reaction's contribution is
third in this layering, not first: it preserves the speed of the verb
against the requirements that would otherwise force I/O into it.

The structural mechanism is, in any instantiation that exercises the
partition, that the deferral artifact runs outside the verb's write
semaphore. In Puppeteer, a Reaction does not run inside the write
semaphore of \texttt{PerformCommand}; it runs after the command has
persisted, on a separate thread (in the case of \texttt{Cue}) or on
demand against the journal (in the case of \texttt{Job}). The verb
returns to the caller without waiting for any deferred step to begin.
The actor's serial-mailbox invariant --- which would otherwise impose a
throughput cap whenever a verb does anything slow --- becomes a
throughput floor: peak throughput is what verbs are \emph{able} to
achieve once the slow work has been hoisted out by I/O hoisting at the
caller, by saga-phasing between events, or by Reactions on the deferred
branch.

What §5.2 names, then, is the third layer of the explanation. The first
two --- in-memory state, rich local methods --- are already present in
any actor framework that respects them. What this paper adds is the
condition under which they survive contact with operational
requirements: a Reaction primitive cheap enough to absorb every effect
that would otherwise be tempted to live inside the verb. Section 6
develops the mechanism in code references; the structural commitment is
summarised in \texttt{Reaction.cs:715-832} and
\texttt{ActorHandler.cs:1662-1733}.

\subsubsection{5.3 Opt-in eventual
consistency}\label{opt-in-eventual-consistency}

The third consequence honors the forward-reference from
\href{01-anti-porosity.md}{Paper 1 §6}: the eventual consistency that
classical CQRS treats as a defining commitment becomes, in actor-native
systems, a shadow of the partition rather than its design choice.

The classical model is opt-out. Read and write stores are physically
separated; their synchronisation is asynchronous by default; the
application accepts the contract, including its observability gap, as
the price of the architecture. Greg Young's \emph{CQRS Documents} make
the trade-off explicit: stronger contracts are recoverable, but only by
application-level effort layered over the framework's defaults.

The actor-native model is opt-in. The actor itself is the read model ---
\texttt{PerformCmd} and \texttt{PerformQry} operate on a single
in-memory state (claim 1) --- so within the actor's boundary, reads and
writes are strongly consistent without any ceremony. The
eventual-consistency contract appears only at the deferred branch: when
the developer declares a Reaction, what happens after the verb's
response is named explicitly. That act of naming is what the contract
is. Each deferred effect is documented at the point of deferral, with a
name and a placement, and the system makes no further claim about its
observability than what the developer has stated.

This inversion is structural, not stylistic. A reader of the system can
list the eventual-consistency windows of any given verb by reading the
Reactions declared near it; the runtime offers nothing else. A reader of
a classical CQRS system needs to know, separately, the topology of the
read store, the synchronisation policy, the lag profile under load, and
the application's strategies for compensating against them. The
Puppeteer reader has, in the worst case, the source of the Reaction; the
classical reader has, in the best case, an architecture diagram and a
runbook.

\subsection{6. Implementation in the Puppeteer
framework}\label{implementation-in-the-puppeteer-framework}

\emph{{[}Skeleton --- reorganized 2026-05-06 to reflect the technical
handoff verified against \texttt{Puppeteer/EventSourcing/Follower/} and
\texttt{ActorHandler.cs}. Path:line citations have been verified against
the current Puppeteer mainline.{]}}

\subsubsection{6.1 The two modes: Cue and
Job}\label{the-two-modes-cue-and-job}

The \texttt{ReactionMode} enum exposes two values, \texttt{Cue} and
\texttt{Job}, declared in \texttt{Reaction.cs:19-23}. The two are not
points on a continuum chosen by the developer at runtime; they are
structurally distinct execution modes, picked when the Reaction is
defined and binding the rest of its semantics.

\texttt{Cue()} is the live-actor mode. The framework registers the
Reaction via \texttt{actorHandler.AddRecordWrittenCallback(...)} after
an initial catch-up batch (\texttt{Reaction.cs:563-568}); from that
point on, every record written to the journal wakes the Reaction's push
loop with the new entry. The loop runs on a separate thread and sleeps
adaptively with backoff between empty checks. Because the Reaction
shares the actor's in-memory state with the director --- the same
rehydration, no separate automaton --- its access is sub-second and
inexpensive while the actor is hot. \texttt{ReadForward()} is mandatory:
a \texttt{Cue} cannot read backward, since its purpose is to react to
the actor's ongoing trajectory, not to its past.

\texttt{Job()} is the journal-driven mode. There is no push loop. The
Reaction runs only when an external scheduler or a call to
\texttt{Reactions.Execute()} triggers it. It reads the journal directly
rather than the actor's memory, which makes the Reaction
\emph{portable}: any process with read access to the replicated journal
can host the work. The case enabling specialised workers ---
\emph{``this pod runs the followers for campaigns; that pod runs the
heavier analytics followers; another pod activates only periodically to
handle older events''} --- is built on this primitive. Both
\texttt{ReadForward()} and \texttt{ReadBackward()} are admissible.

\texttt{Job} does not rehydrate the actor. When the Reaction needs
domain-computed data, the canonical mechanism is \texttt{expose} (§6.7)
--- the command pre-writes the data to the journal alongside the action,
and the Reaction reads it during its own rehydration without ever waking
the actor. This keeps \texttt{Job} cheap to host on a journal-only
follower; reaching into the live actor would defeat the cost model that
justifies the mode in the first place. The non-rehydration boundary is
therefore not a missing feature but a structural commitment that
preserves the distinction between the two modes.

The two modes correspond to two structurally distinct points on the
proximity/urgency axis introduced in claim 11. A comparative table:

{\def\LTcaptype{none} % do not increment counter
\begin{longtable}[]{@{}
  >{\raggedright\arraybackslash}p{(\linewidth - 4\tabcolsep) * \real{0.3333}}
  >{\raggedright\arraybackslash}p{(\linewidth - 4\tabcolsep) * \real{0.3333}}
  >{\raggedright\arraybackslash}p{(\linewidth - 4\tabcolsep) * \real{0.3333}}@{}}
\toprule\noalign{}
\begin{minipage}[b]{\linewidth}\raggedright
Aspect
\end{minipage} & \begin{minipage}[b]{\linewidth}\raggedright
\texttt{Cue}
\end{minipage} & \begin{minipage}[b]{\linewidth}\raggedright
\texttt{Job}
\end{minipage} \\
\midrule\noalign{}
\endhead
\bottomrule\noalign{}
\endlastfoot
State source & Actor's live memory (same process) & Rehydration from
journal \\
Lifts a separate automaton & No & Yes (follower hydrates) \\
Trigger & \texttt{OnRecordWritten} push callback & Explicit
invocation \\
Latency & Sub-second & Batch / on-demand \\
Distributable & No & Yes \\
Direction & \texttt{ReadForward} only & Forward or backward \\
\end{longtable}
}

The choice between them is an architectural decision, not a tuning knob.
A Reaction modeled as \texttt{Cue} cannot transparently become a
\texttt{Job}: the two assume different things about what is in memory,
who hosts the work, and whether the actor must be live for the Reaction
to make progress. This is the §2.3 axis materialised --- partition and
placement are one act because, at the moment of partitioning, the
developer is also choosing which of these two modes will carry the
deferred work. Both modes operate as journal-indexed cursors per §3.6:
the choice between them governs \emph{who walks the journal and when},
not whether the cursor exists.

\subsubsection{6.2 Static pattern matching against the domain
libraries}\label{static-pattern-matching-against-the-domain-libraries}

The Reaction's binding to the domain is established at build time, not
at runtime. Three layers of validation make this concrete.

The first layer is static parsing in \texttt{Pattern.cs}. Every
reference of the form \texttt{{[}\_:Class{]}}, \texttt{{[}var:Class{]}},
or \texttt{Class(param:Type,\ ...)} is resolved against the actor's
\texttt{Libraries}, case-insensitively. If the class or the constructor
signature does not exist in the domain library, the Reaction fails to
construct --- there is no runtime path that can recover from a missing
type. The Libraries here are the same domain types that
\href{01-anti-porosity.md}{Paper 1} named as the legitimate tool inside
the library and that \href{02-program-value-separability.md}{Paper 2}
treated as the substrate against which DSL programs name their
operations; the Reaction is one more language layer that compiles
against that substrate.

The second layer runs per event during rehydration.
\texttt{Reaction.SolveActionReferences()} (\texttt{Reaction.cs:997})
parses the script of each \texttt{ActionEventData} exactly once and
caches it under an LRU of 100 entries; identifiers that the pattern
captured are marked \texttt{IsParameter=true} so subsequent stages can
reuse them without re-parsing. The cache amortises the cost of parsing
the same script many times across many events; the LRU bound prevents
memory growth in workloads with high script diversity.

The third layer is the runtime walk.
\texttt{MatchTree.TryMatchAtLevel()} traverses the tree of stages by BFS
(\texttt{WithSharedHydration}, where captures from earlier stages remain
visible to later ones) or by DFS (\texttt{WithIndependentHydration},
where each branch carries its own hydration context). The choice between
the two is part of the Reaction's definition and shapes how identifier
propagation behaves across \texttt{Seek}, \texttt{ThenSeek}, and
\texttt{ThenFinalSeek} stages.

The parser, importantly, does not execute the script's expressions. It
only knows the types. Matches over non-literal values are therefore by
structure and type, not by value equality --- with two exceptions:
literals (a constant in the pattern) and \texttt{Id} parameters with
\texttt{IsParameter==true} (the captured identifiers of §6.3). The
composition rules round out the layer: a Reaction has at least one
\texttt{Seek} and any number of \texttt{ThenSeek} stages, optionally
terminated by a \texttt{ThenFinalSeek}; each \texttt{Seek} may carry
several \texttt{OnMatch} clauses, all of which must match against the
same underlying script. The build-time validation stops a Reaction from
being defined against a domain it does not fit; the runtime walk only
ever evaluates patterns the Libraries have already approved.

\subsubsection{6.3 Multi-event patterns: trajectory, not
event}\label{multi-event-patterns-trajectory-not-event}

The simplest non-trivial Reaction names three stages of an order's life
and correlates them by both the order instance and a primitive
identifier:

\begin{Shaded}
\begin{Highlighting}[]
\NormalTok{actor}\OperatorTok{.}\FunctionTok{Reactions}
    \OperatorTok{.}\FunctionTok{DefineReaction}\OperatorTok{(}\StringTok{"OrderCompleted"}\OperatorTok{)}
    \OperatorTok{.}\FunctionTok{Job}\OperatorTok{().}\FunctionTok{DirectorOnly}\OperatorTok{().}\FunctionTok{ReadForward}\OperatorTok{().}\FunctionTok{WithSharedHydration}\OperatorTok{()}
    \OperatorTok{.}\FunctionTok{Seek}\OperatorTok{(}\StringTok{"Initiate"}\OperatorTok{)}
        \OperatorTok{.}\FunctionTok{OnMatch}\OperatorTok{(}\StringTok{"[\_:Customer].InitiateOrder(order:Order, number:int)"}\OperatorTok{)}
    \OperatorTok{.}\FunctionTok{ThenSeek}\OperatorTok{(}\StringTok{"Confirm"}\OperatorTok{)}
        \OperatorTok{.}\FunctionTok{OnMatch}\OperatorTok{(}\StringTok{"order.Confirm(number)"}\OperatorTok{)}
    \OperatorTok{.}\FunctionTok{ThenFinalSeek}\OperatorTok{(}\StringTok{"Pay"}\OperatorTok{)}
        \OperatorTok{.}\FunctionTok{OnMatch}\OperatorTok{(}\StringTok{"order.Pay(number, [\_:PaymentMethod])"}\OperatorTok{)}
    \OperatorTok{.}\FunctionTok{Program}\OperatorTok{.}\FunctionTok{Emit}\OperatorTok{(}\StringTok{"..."}\OperatorTok{);}
\end{Highlighting}
\end{Shaded}

Read aloud, the pattern says: \emph{``when, in the actor's history, a
customer initiated an order with a given number N, and that same order
with that same number N was later confirmed and paid --- react.''}

Five details of this small example carry the weight of the primitive.

\emph{Identifier propagation.} The captured \texttt{order} named at
\texttt{Seek("Initiate")} is the same \texttt{order} \emph{instance}
that must appear at \texttt{ThenSeek("Confirm")} and at
\texttt{ThenFinalSeek("Pay")}; the captured \texttt{number} is the same
primitive \emph{value}. The matcher correlates by typed identity for
object captures and by literal equality for primitive captures --- both
kinds of identifiers, once captured, are treated as fixed for the
remainder of the trajectory. A \texttt{Confirm} invoked with a different
number, or for a different order instance, would not advance the
trajectory; the trajectory is locked to \emph{this} order with
\emph{this} number once the first stage matches.

\emph{Non-contiguity.} Between \texttt{Confirm} and \texttt{Pay} the
actor may have processed any number of unrelated commands; the matcher
steps over them. The trajectory is a path through the journal, not a
contiguous window.

\emph{One match per stage.} Each \texttt{Seek} and \texttt{ThenSeek}
matches exactly once. For one-to-many correlations --- an order with
several items being added before payment --- the framework uses
\texttt{RepeatSeek}, which is restricted to the initial position of the
pattern. The surrounding language and its limits are documented in §6.4.

\emph{Pattern over conduct.} What the matcher inspects is the sequence
of invocations the actor recorded --- its conduct as a sequence of
\texttt{PerformCommand} calls --- not the body of any one message. The
Reaction sees what the actor \emph{did}, not what arrived in its
mailbox.

\emph{Static binding.} \texttt{Customer}, \texttt{Order},
\texttt{PaymentMethod} are types from the actor's domain library,
validated at build time against \texttt{Libraries}; the symbol table of
the actor is available to the matcher, so polymorphism and type
relations are respected without further configuration.

Three properties together --- \emph{trajectory}, \emph{static binding},
\emph{typed propagation} --- make this primitive qualitatively distinct
from event handlers, grain reminders, and CEP rules over external
streams. The first two §3.5 already named; the third is what makes the
example above readable as a saga rather than as three loosely-coupled
callbacks.

\subsubsection{6.4 Filters, time,
quantifiers}\label{filters-time-quantifiers}

The matcher accepts three families of refinement on top of the base
pattern.

\emph{Where clauses}. \texttt{Where(expression)}
(\texttt{ReactionEngine.cs:47-60}) chains a boolean condition onto a
\texttt{Seek}. Inside the expression the developer has access to
\texttt{@Now} (the event's timestamp, non-nullable), \texttt{@EntryId}
(the journal entry's identifier, non-nullable), the identifiers captured
by the pattern, and the global \texttt{time} instance of the framework's
\texttt{Temporal} type. \texttt{time.Days(n)}, \texttt{time.Hours(n)},
\texttt{time.Minutes(n)} and so on (\texttt{Reaction.cs:525}) construct
\texttt{TimeSpan} values for relative windows. The framework rejects
always-false comparisons at build time: comparing \texttt{@Now} or
\texttt{@EntryId} against \texttt{null} raises a
\texttt{LanguageException}, since their types do not admit a null
sentinel.

\emph{Quantification}. The base pattern matches each \texttt{Seek}
exactly once; for one-to-many correlations the framework provides
\texttt{RepeatSeek(name)} (\texttt{Reaction.cs:318-329}), which
accumulates multiple matches at a single level before transitioning to
the next stage. The chained modifiers
(\texttt{ReactionEngine.cs:117-149}) are \texttt{.GroupBy(varName)}
(group accumulated matches by the value of a captured variable),
\texttt{.Accumulate(name)} (expose the accumulated matches under a name
visible to the next stage), \texttt{.Distinct(varName)} (filter
duplicates by a captured variable), and \texttt{.Until(...)} taking
either an integer count or a \texttt{TimeSpan}. The constraint is
positional and intentional: \texttt{RepeatSeek} is admissible only at
the start of the pattern. There is no \texttt{ThenRepeatSeek} mid-chain
--- a topology of \emph{one event, then several, then one} must either
model the repetition first or decompose into two cooperating Reactions.

\emph{Lifecycle}. \texttt{SetExpirationDate(DateTime)}
(\texttt{Reaction.cs:619-623}) shuts down the Reaction wholesale after a
date; \texttt{IsExpired} (\texttt{Reaction.cs:627-634}) returns
\texttt{true} once \texttt{DateTime.UtcNow} exceeds the expiration.
\texttt{SetActive(bool)} (\texttt{Reaction.cs:613-617}) is the manual
override. These two differ in scope from \texttt{Until(timeout)}:
\texttt{Until} closes a single accumulation window inside a
\texttt{RepeatSeek}; \texttt{SetExpirationDate} and \texttt{SetActive}
toggle the entire Reaction. The two scopes are useful for different
problems --- \texttt{Until} lives in the pattern;
\texttt{SetExpirationDate} lives outside it.

\subsubsection{6.5 The three action
planes}\label{the-three-action-planes}

A Reaction takes exactly one action when its pattern matches. That
action is expressed through one of three named \emph{planes}, each
describing a different surface of the system the matched trajectory is
permitted to affect.

{\def\LTcaptype{none} % do not increment counter
\begin{longtable}[]{@{}
  >{\raggedright\arraybackslash}p{(\linewidth - 4\tabcolsep) * \real{0.3333}}
  >{\raggedright\arraybackslash}p{(\linewidth - 4\tabcolsep) * \real{0.3333}}
  >{\raggedright\arraybackslash}p{(\linewidth - 4\tabcolsep) * \real{0.3333}}@{}}
\toprule\noalign{}
\begin{minipage}[b]{\linewidth}\raggedright
Plane
\end{minipage} & \begin{minipage}[b]{\linewidth}\raggedright
What the verb touches
\end{minipage} & \begin{minipage}[b]{\linewidth}\raggedright
What it cannot touch
\end{minipage} \\
\midrule\noalign{}
\endhead
\bottomrule\noalign{}
\endlastfoot
\texttt{Program} & the actor's domain libraries, read-only & the
journal, the actor's state, any global \\
\texttt{Causation} & the actor's own journal (a journaled cross-actor
send) & the domain libraries or the metadata plane \\
\texttt{Metadata} & the journal's elision register & the domain state or
any external effect \\
\end{longtable}
}

A Reaction is therefore not an unbounded scripting surface attached to a
pattern. It is a constrained declaration of which surface of the system
the matched trajectory is permitted to affect. The builder enforces that
only one plane can be configured per Reaction; configuring a second
plane after one has already been set is a construction error. The
\emph{exactly-one-action} rule is structural, not conventional --- the
developer cannot accidentally route a single match into two different
planes because the surface refuses the second configuration call at
build time.

\paragraph{Program --- read-only execution against the domain
libraries}\label{program-read-only-execution-against-the-domain-libraries}

The \texttt{Program} plane executes a script under the same regime as a
query. The parser runs in query mode, the script is forbidden from
declaring globals, the \texttt{expose} keyword is rejected, the journal
is not written, and the runtime takes a read lock that runs in parallel
with other queries and emits.

The script may invoke arbitrary external technology --- an HTTP client,
a Kafka producer, a SignalR hub, a BI sink. What the language denies is
any path by which those calls could feed back into the actor's state.
This is not a discipline rule kept by convention; it is rejected at
parse time, before the developer has the chance to write the leak. A
Reaction may observe the actor and affect the outside world, but it
cannot feed that effect back into the actor.

The optional \texttt{when:} guard performs a second read-only check at
firing time:

\begin{itemize}
\tightlist
\item
  the pattern asserts that \emph{this trajectory occurred in history},
\item
  the guard asserts that \emph{it still makes sense to react now}.
\end{itemize}

The distinction matters under replay, catch-up after downtime, and
replica rehydration --- the three situations in which the temporal
distance between history and the present can be arbitrary. The pattern
is bound to the trajectory; the guard is bound to the clock.

A worked example. Consider a Reaction that, after a confirmed purchase,
pushes a real-time toast notification to the customer's UI:

\begin{Shaded}
\begin{Highlighting}[]
\NormalTok{seller.Reactions.DefineReaction("NotifyPurchaseToast")}
\NormalTok{    .Cue().Company().ReadForward()}
\NormalTok{    .Seek("Purchase")}
\NormalTok{        .OnMatch("[s:Seller].purchase($orderId, $date, $amount, $customer)")}
\NormalTok{    .Program.Emit(}
\NormalTok{        "notifier.Push(@customer, @orderId);",}
\NormalTok{        when: "check (notifier.IsFresh(@date)) WARNING \textquotesingle{}toast stale, drop\textquotesingle{};"}
\NormalTok{    );}
\end{Highlighting}
\end{Shaded}

The domain verb \texttt{Seller.purchase(...)} carries no reference to
the notification hub, to UI presence, to toast timing, or to any
operational concern of the notification layer. All of that lives inside
the Reaction: the call to \texttt{notifier.Push(...)} invokes the
extrinsic technology, and the guard decides at firing time whether the
toast still has an audience. Under the steady path the push fires within
seconds of the journal entry; under catch-up or replay, the same match
arrives at the matcher hours or days later, the guard returns false, and
the emit is suppressed. The freshness policy is plain code in a library
class --- \texttt{ToastNotifier} --- that the actor sees through its
\texttt{Libraries} but that the domain verb never references. The
segregation is guaranteed because there is no syntactic route by which
this Reaction could import the notifier back into
\texttt{Seller.purchase}.

\paragraph{Causation --- journaled cross-actor
continuation}\label{causation-journaled-cross-actor-continuation}

The \texttt{Causation} plane is the journaled cross-actor send. A
Reaction on this plane writes exactly one sentence into the sender's own
journal instructing another actor; the sender records the causation, the
receiver processes the message independently through its own endpoint.
Cross-actor coordination therefore appears in the program as a domain
fact rather than as infrastructure --- neither a saga step nor a tracing
span nor a choreography event, but a sentence in the actor's history.

This plane is developed in detail in
\href{04-cross-actor-continuity.md}{Paper 4}, which introduces the
primitive, names the structural conditions under which it preserves
cross-actor program continuity, and presents the comparative case study
against sagas, choreography, and distributed tracing. Here it is named
only as the second of the three surfaces a Reaction may touch.

\paragraph{Metadata --- declaring a trajectory
closed}\label{metadata-declaring-a-trajectory-closed}

The \texttt{Metadata} plane marks the events covered by the pattern as
elided. After a trajectory has been correlated, its component events no
longer participate in any future decision; subsequent rehydrations walk
past them.

This is the conceptual bridge to the \emph{skips} introduced in
\href{01-anti-porosity.md}{Paper 1}. What appeared there as a
journal-density optimisation is here the natural outcome of declarative
correlation. The developer is not cleaning the journal after the fact.
The developer is declaring \emph{``this trajectory is closed; its
component events no longer decide anything.''} The skip is the
consequence of that declaration, not its cause.

\paragraph{Parameter injection}\label{parameter-injection}

\texttt{WithParameters(...)} allows captures from the pattern matching
to be modified or augmented before the action runs, regardless of which
plane is configured. The parameters are pre-populated with the pattern's
captures, so the idiomatic use is to enrich them rather than to
construct them from scratch.

\paragraph{Why the planes matter}\label{why-the-planes-matter}

Without named planes, a Reaction would be an unbounded script runner
attached to history --- a feature whose semantics would depend on which
calls the developer happened to write. With named planes, a Reaction
becomes a one-sentence statement about which layer of the system is
allowed to change when a historical pattern is recognised. The same
primitive that detects the trajectory also declares the surface of
effect, and the surface of effect is selected from a closed set of three
rather than from the open set of all calls a script could make.

That restriction is what makes the segregation between

\begin{itemize}
\tightlist
\item
  domain verbs,
\item
  cross-actor causation,
\item
  operational side-effects, and
\item
  journal maintenance
\end{itemize}

reliable across teams and across time. The reliability is not a property
of the developer's discipline. It is a property of the surface the
language permits the developer to address.

\subsubsection{6.6 Activation: where the Reaction
runs}\label{activation-where-the-reaction-runs}

After \texttt{Cue()} or \texttt{Job()}, the developer chooses an
activation scope (\texttt{Reaction.cs:67-80}). \texttt{DirectorOnly()}
runs the Reaction only on the primary replica --- appropriate when the
work must be unique across the topology, such as orchestrations that
should not be duplicated. \texttt{CastOnly()} runs only on followers
(secondary replicas), enabling the specialised-workers case named in
§6.1: followers that hydrate against the journal alone, hosting
\texttt{Job} Reactions that the director does not need to evaluate.
\texttt{Company()} runs on the director and on all followers;
appropriate when each replica should hold its own evaluation of the
trajectory, perhaps for local cache invalidation or for redundant emit
paths.

Activation composes with the StageManager --- the topology director of
the Choreography module. The StageManager decides who is director and
who is follower at any moment; the Reaction inherits that identity to
decide whether to fire. A Reaction defined \texttt{CastOnly} becomes
silent on the director and active on whatever process the StageManager
has elected as a follower; a director-to-follower handoff therefore
changes the population of running Reactions without any further action
by the developer. The deployment-time mechanism that keeps such handoffs
safe is \texttt{Theater.aliveGate} (claim 8); it is treated in detail in
\href{05-zero-downtime.md}{Paper 5}.

\subsubsection{\texorpdfstring{6.7 \texttt{expose} versus
\texttt{Program.Emit}}{6.7 expose versus Program.Emit}}\label{expose-versus-program.emit}

The two mechanisms are easy to confuse and are not the same thing. They
are addressed in the same direction --- feeding downstream systems data
the actor has computed, without forcing those systems to rehydrate the
actor --- but from opposite sides of the journal.

\texttt{expose} is a keyword of the actor's language, valid only inside
a \texttt{PerformCommand}. When the verb runs, the script can mark a
value as exposed, and the framework persists it into the
\texttt{ExposeData} of the resulting journal entry alongside the
\texttt{Action}. It is the \emph{writer's} externalisation point: the
cashier confirming an order can expose the order number, the calculated
total, the resolved campaign --- and downstream systems can consume
those values directly from the journal without having to walk the
actor's state.

\texttt{Program.Emit} is the Reaction's \emph{reader's} externalisation
point. A Reaction running over the journal --- whether \texttt{Cue} or
\texttt{Job} --- has access to the \texttt{ExposeData} of the events it
traverses (\texttt{Reaction.cs:573}, \texttt{includeExposeData=true}),
and it can use those values in its \texttt{Where} clauses, in its
captures, and in the parameters of its emit. What it cannot do is
\emph{produce} an \texttt{expose}: \texttt{PerformEmit} blocks the
keyword by parsing with \texttt{isQuery:true} (§6.5). The asymmetry is
structural and intentional. \texttt{expose} is the writer's hook into
the journal; \texttt{Program.Emit} is the reader's hook out of it.

\subsubsection{6.8 What does not exist --- honest
limits}\label{what-does-not-exist-honest-limits}

A technical paper that reports what a primitive supports without
reporting what it does not is incomplete. The Reaction primitive does
not, today, provide:

\begin{itemize}
\tightlist
\item
  \textbf{Negative or absence patterns.} No \texttt{Without},
  \texttt{NotMatch}, \texttt{Negate}, \texttt{Unless}, \texttt{Missing}.
  \emph{``X happened but Y did not happen afterwards''} is not directly
  expressible. The conventional workaround is a periodic \texttt{Job}
  that matches \texttt{X} and consults the actor's state via
  \texttt{Program.Emit} with a \texttt{when:} check to confirm
  \texttt{Y} has not occurred.
\item
  \textbf{Traditional regex-style quantifiers.} No \texttt{AtLeast(n)},
  \texttt{AtMost(n)}, \texttt{Exactly(n)}, \texttt{OneOrMore},
  \texttt{Optional}. The only quantification primitive is
  \texttt{RepeatSeek} with \texttt{Until(count)} or
  \texttt{Until(timeout)}, restricted to initial position.
\item
  \textbf{Absolute time windows} as a dedicated method
  (\texttt{Within(start,\ end)}). Expressed instead via
  \texttt{Where("@Now\ \textgreater{}=\ ...\ \&\&\ @Now\ \textless{}=\ ...")}.
\item
  \textbf{Optional \texttt{ThenSeek}.} A \texttt{ThenSeek} cannot be
  marked optional. Optional branches are written as separate Reactions.
\item
  \textbf{\texttt{Job}-mode rehydration of the actor.} A \texttt{Job}
  Reaction reads only the journal; it does not wake or rehydrate the
  actor itself. When domain-computed state is needed by the Reaction,
  the canonical mechanism is \texttt{expose} (§6.7), which pre-writes
  the data alongside the command. Silent rehydration inside \texttt{Job}
  is intentionally absent --- collapsing this boundary would erase the
  structural distinction between \texttt{Cue} and \texttt{Job} (§6.1)
  and reintroduce the cost the design was avoiding.
\item
  \textbf{Cross-actor dispatch with journal-recorded causation ---
  historically a gap, now resolved.} When v0.1 of this paper was
  published, this was a recognized limitation: cross-actor causation was
  mediated by \texttt{ITransport} infrastructure outside the sender's
  journal, leaving the broker as the seam. Since v0.1, the framework has
  gained the third action plane --- \texttt{Causation.Continue} (the
  Tell primitive) --- that records the cross-actor dispatch as a
  sentence in the sender's own journal while leaving the transport
  choice to the developer's \texttt{ITransport} and the target actor's
  processing entirely to its own endpoint. The construct, the design
  rationale, and the comparison against sagas, choreography, and
  distributed tracing are the subject of
  \href{04-cross-actor-continuity.md}{Paper 4}. The bullet is preserved
  here as historical record of how the gap was named before being
  filled; the present paper does not develop the primitive (it is
  referenced from §6.5 as the third plane).
\end{itemize}

These absences are honest design boundaries of the present
implementation. Some are gaps in the original sense --- pending features
that may arrive --- and some, like the last one, are deliberate
commitments without which the surrounding primitive would lose
coherence. The framework's scope is what the framework supports; what it
does not support is what the developer must build over it or model
differently.

\subsection{7. Comparison with Akka and
Orleans}\label{comparison-with-akka-and-orleans}

\subsubsection{7.1 Akka actors and
Become/Receive}\label{akka-actors-and-becomereceive}

Akka's actors model behavior change with \texttt{Become} and
\texttt{Unbecome}, swapping the receive partial function as the actor
moves between modes. A booking actor might \texttt{Become(opened)}, then
\texttt{Become(confirmed)}, then \texttt{Become(paid)}; each receive
defines a different set of messages the actor accepts and how it handles
them. The partition this expresses is across \emph{actor states} ---
what the actor will accept next changes with the state --- and the
developer's mental model is a state machine the actor walks through.

Reactions express a different partition. The actor's
\texttt{PerformCommand} is one verb, not many --- it does not
\texttt{Become} between calls --- and the partition that Reactions add
is across the \emph{event boundary}: when a command has been processed
and persisted, what happens afterwards is the Reaction's domain. Akka's
mechanism is suited to actors whose interface changes as state evolves;
Puppeteer's mechanism is suited to actors whose interface is stable but
whose effects after each verb are rich. Neither subsumes the other; they
answer different questions.

\subsubsection{7.2 Orleans grains, timers, and
reminders}\label{orleans-grains-timers-and-reminders}

Orleans' grains support timers (in-process, not durable across grain
reactivation) and reminders (durable, scheduled). Both share the
placement intent of \texttt{Job}: deferred work, possibly on another
silo, scheduled rather than co-located. The differentiator is friction.
Registering an Orleans timer or reminder is procedural --- a method
call, a callback signature, a registration step in the grain's
activation lifecycle --- and the work itself is general C\# code. A
Reaction is a few lines of fluent DSL declared next to the verb that
triggers it (§4), with a static pattern that compiles against the domain
library (§6.2) and a constrained set of actions (§6.5). The framework's
contribution is not a feature Orleans lacks but a friction profile
Orleans does not aim for; the developer who wants the partition
exercised by default in Orleans will write more code than the developer
who wants the same in Puppeteer.

\subsubsection{7.3 Where the design spaces converge and
diverge}\label{where-the-design-spaces-converge-and-diverge}

Akka, Orleans, and Puppeteer address the same structural pressures: the
actor's serial mailbox, the cost of state hydration, the cost of
distributing work across nodes, the difficulty of keeping fault
isolation honest. The three answer those pressures with different design
commitments. Akka commits to a flexible behavior model in which the
actor's interface itself can evolve; Orleans commits to virtual actors
that materialise on demand, with placement and scheduling as runtime
concerns. Puppeteer commits to an externalised DSL with a homoiconic
journal (\href{02-program-value-separability.md}{Paper 2}) and an
anti-porous domain (\href{01-anti-porosity.md}{Paper 1}) --- and
Reactions are the language layer of that commitment: declarative pattern
matching over the actor's trajectory (§3.5), with a strict separation
between the actor's command path and what runs afterwards (§5.2). The
Puppeteer position is one point in a non-trivial design space, not its
peak.

Nothing in the present paper claims Reactions are the only realization
of the construct. An Akka or Orleans extension that exposed a
declarative DSL for trajectory matching with build-time binding to the
host language's types would satisfy the same construct; this paper
presents one realization, not the realization.

\subsection{8. Counter-arguments}\label{counter-arguments}

\subsubsection{8.1 ``Eventual consistency is a known anti-pattern in
business
systems''}\label{eventual-consistency-is-a-known-anti-pattern-in-business-systems}

The objection treats eventual consistency as a uniform property of the
system, taken or rejected wholesale, and reads the warning literature on
it as applying everywhere it appears. The framing is correct for
\emph{opt-out} eventual consistency, the kind classical CQRS introduces
by separating read and write stores: every read is potentially stale,
every write is potentially unobserved, and the application has to
recover stronger contracts wherever the business demands them. Under
that framing, the warning is well-earned.

The model developed in this paper inverts the framing (§5.3). Reads and
writes inside the actor are strongly consistent --- \texttt{PerformCmd}
and \texttt{PerformQry} operate on the same in-memory state. Eventual
consistency appears only at the deferred branch, named by the developer
one deferral artifact at a time, with a placement attached. There is no
system-wide consistency contract that the application must compensate
for; there is, instead, a list of named windows the developer has chosen
to admit. The literature warning applies to the framework that imposes
eventual consistency as a default; a reader of an instantiation that
exercises this construct can answer ``what is eventually consistent
here?'' by reading the deferral artifacts declared near the verb and
nothing else.

The honest concession (claim 6's 20\%): some business domains require
that every observable effect of a verb commit in lockstep with the
response --- banking core ledgers, certain regulated transaction
systems. The opt-in framing does not soften this; in those domains, the
developer cannot extend a license that does not exist, and the
partition's primitive is not the right primitive (§When NOT, A). The
objection lands where it lands; it does not generalise to every business
system that adopts the framework.

\subsubsection{8.2 ``Reactions are async events with a different
name''}\label{reactions-are-async-events-with-a-different-name}

The objection holds only if Reactions are read at their surface --- as
something that runs after a verb, on a different thread, with a callback
shape. At that surface they look like async events, and the objection
follows.

Read at the structural level of the construct, the artifact realizing
the partition is something else. An async event handler responds to a
single message; the construct's matching primitive matches a
\emph{trajectory} --- a sequence of one or more stages naming a saga, a
lifecycle, or an anomaly through the actor's history (claim 10, §6.3).
The pattern is parsed against the domain's libraries at build time, not
against text at runtime: a reference to a class the domain does not
contain fails to construct (§6.2). Identifiers captured in one stage are
typed and propagated to the next, which is what correlates the events
into a story --- an async event handler has no analogue. The placement
axis with two canonical points (in Puppeteer's instantiation,
\texttt{Cue} and \texttt{Job}) carries activation scopes
(\texttt{DirectorOnly} / \texttt{CastOnly} / \texttt{Company}) that
compose with the StageManager's topology (§6.6); an async event has none
of that surface. The action plane choices (\texttt{Metadata.Elide},
read-only \texttt{Program.Emit} with optional \texttt{when:} check, and
the journaled cross-actor \texttt{Causation.Continue}) are constrained
primitives --- particularly the read-only \texttt{Program.Emit} whose
underlying \texttt{PerformEmit} structurally blocks \texttt{expose} and
journal writes (§6.5) --- designed to enforce the categorical
segregation of §3.3 in the parse phase, not at runtime.

The honest concession (claim 6's 20\%): the surface mechanism --- work
happens after the verb returns, on a separate thread or on demand --- is
shared with async event systems. A reader who looks only at the trigger
sees something familiar. The argument of this paper is that the trigger
is the least interesting part of the primitive.

\subsubsection{8.3 ``Even with low-friction Reactions, the developer can
pollute the
domain''}\label{even-with-low-friction-reactions-the-developer-can-pollute-the-domain}

The objection is correct in spirit: friction reduction is not a
substitute for discipline. A determined developer can put
\texttt{sendEmail(...)} inside a \texttt{PerformCommand} even when a
Reaction one line below would do the job; the framework cannot prevent
it. Were the argument that friction reduction \emph{replaces} developer
judgment, the objection would land.

The argument is different. Friction reduction is the lever the framework
can pull; developer judgment is the lever it cannot. When friction is
high --- when the right thing costs more than the wrong thing --- no
amount of training, code review, or architectural pep-talk holds at
scale. When friction is low, training and review have something to lean
on; the developer who chooses inline anyway is doing so against an
obvious cheaper alternative, and the choice becomes legible to reviewers
as a deliberate departure rather than a path of least resistance.

The honest concession is at the seam where friction reduction reaches
its limit: a careless team will write polluted code under any framework,
and a disciplined team will write clean code under any framework. The
framework's contribution is the \emph{gradient} between these --- making
clean code the default, pollution the exception, and the cost of getting
it wrong recoverable rather than irreversible. The objection becomes an
argument for friction reduction, not against it: in a world where
developer discipline is unreliable, the only thing the architecture can
do is shape the cost surface so that the path of least resistance is
also the path of least damage.

\subsubsection{8.4 ``DSL scripts repeat lines across endpoints --- this
violates
DRY''}\label{dsl-scripts-repeat-lines-across-endpoints-this-violates-dry}

The objection assumes DRY as a textual prohibition. In its original
formulation (Hunt \& Thomas, 1999), DRY prohibits the duplicate
representation of \emph{knowledge}, not the appearance of similar lines:
\emph{``every piece of knowledge must have a single, unambiguous,
authoritative representation within a system.''} The domain library
lives once, in OOP; DSL scripts are the language in which an actor is
invoked, structurally analogous to SQL. Two SQL queries that share JOIN
patterns are not in violation of DRY --- they are compositions over a
shared schema. The same is true of two DSL endpoints that share
invocation lines: the knowledge is in the library that the lines invoke,
not in the lines themselves.

Concession (the 20\%): if two DSL endpoints share \emph{domain logic}
--- not invocation patterns but actual computation --- that logic should
refactor into a method of the library. The line is between invocation
and implementation, not between unique and repeated text.

\subsubsection{8.5 ``Reactions should be able to rehydrate the actor for
richer state
queries''}\label{reactions-should-be-able-to-rehydrate-the-actor-for-richer-state-queries}

The choice between modes is the structural distinction itself --- full
state queries belong to the live-actor mode, not to the journal-driven
one (§6.1). When the data is anticipated by the developer, the journal
pre-write mechanism (\texttt{expose} in Puppeteer, §6.7) puts the data
where the deferral artifact will consume it during its own rehydration,
without ever waking the actor. The journal-only-without-rehydration
boundary is what makes the journal-driven mode cheap to host on a remote
follower; lifting it would collapse the two modes into one fuzzy mode
and reintroduce the cost the original design was trying to avoid.

Concession: there are rare cases of data that no one anticipated would
be needed by a future Reaction, and that already exist in many past
events without an \texttt{expose}. The workaround is operational --- add
\texttt{expose} now, run a migration that replays old events to write
the \texttt{expose} retroactively, or build an auxiliary \texttt{Job}
that materialises a side projection. Costly, but solvable, and rare
enough that it does not justify reopening the primitive.

\subsection{9. Related work}\label{related-work}

\subsubsection{9.1 Actor Model and CQRS
foundations}\label{actor-model-and-cqrs-foundations}

The actor model originates with Hewitt (1973), which introduces the
actor as the universal unit of computation: an entity with private
state, a mailbox for incoming messages, and a single thread of behavior
responding to them. The model has accumulated decades of refinement ---
Agha (1986), Hewitt (2010), Bonér et al.'s \emph{Reactive Manifesto}
(2014) --- but the core remains: state isolated, computation
message-driven, concurrency through the multiplication of actors rather
than the sharing of state. Puppeteer's actor sits in this lineage, with
the addition that the journal of received messages is itself the
persisted artifact rather than a derived log.

Command Query Responsibility Segregation (CQRS) is named in Young
(2010), with intellectual antecedents in Meyer's command-query
separation (1988) and Fowler's reflections on enterprise architecture
(2002). The classical formulation separates write and read stores at the
level of the system architecture, with eventual consistency between them
as the design contract. This paper engages CQRS not as the primary frame
but as the literature that names the consequences of the partition this
paper develops; the ``logical CQRS'' of claim 1 is the classical pattern
read at a different level --- one in-memory state, two verbs over it,
the eventual-consistency contract appearing only at the deferred branch
(§5.3, §8.1).

\subsubsection{9.2 DDD and the anti-pattern
literature}\label{ddd-and-the-anti-pattern-literature}

Domain-driven design originates with Evans (2003), where the
\emph{anemic domain} and \emph{smart UI} antipatterns are diagnosed:
domain types reduced to data containers with behavior emigrated to
service layers, and business rules pulled into the
presentation/transport edge. Vernon (2013) catalogs the implementational
pitfalls of DDD-by-vocabulary-only --- bounded contexts and aggregates
named but never structurally enforced. Fowler (2003, in
\emph{AnemicDomainModel}) gives the antipattern its most cited
articulation. These three sources name the industrial reality §5.1 of
this paper describes: in most enterprises, ``the domain'' does not exist
as a library --- it exists as scattered classes whose names suggest
cohesion and whose substance does not.

Fowler's \emph{Patterns of Enterprise Application Architecture} (2002)
is the relevant catalog for the DTO antipattern: a Data Transfer Object
intended for transport that becomes load-bearing in the domain,
contaminating both. The same volume names the \emph{Distributed Object}
antipattern --- the chatter problem this paper alludes to in §3.1 when
it argues for the SQL-analogous use of a DSL as invocation rather than
transcription.

Hunt and Thomas (1999) define DRY in \emph{The Pragmatic Programmer} as
the prohibition of duplicate \emph{representation of knowledge}, not of
repeated text --- the formulation §3.1 and §8.4 recover against a more
recent industrial reading. Sandi Metz's \emph{AHA --- Avoid Hasty
Abstractions} (2016) provides the natural complement: the cure for
repeated text is sometimes more text, not the wrong abstraction.

\subsubsection{9.3 CEP and adjacent
runtimes}\label{cep-and-adjacent-runtimes}

Complex event processing has a substantial literature --- Luckham's
\emph{The Power of Events} (2002), Etzion and Niblett's \emph{Event
Processing in Action} (2010) --- concerned with detecting patterns over
streams of events, often from heterogeneous sources, often without the
type system of the consuming application available. The pattern matching
machinery of §6 shares structural intent with CEP rules but operates
inside the actor's own type universe: patterns compile against the
domain's \texttt{Libraries} at build time (§6.2), captured identifiers
are typed and propagated (§6.3), and the events being matched are the
actor's own conduct rather than external streams. The closest analogue
from CEP is the typed event pattern of engines that integrate with
strongly-typed host languages; the operative difference is that those
engines treat the pattern as a runtime artifact, while Puppeteer's
pattern is a compile-time artifact that fails to construct against a
domain it does not fit.

Akka and Orleans are the most relevant adjacent runtimes (§7). Their
documentation is the canonical reference for Akka's
\texttt{Become}/\texttt{Receive} and Orleans' grain timers, reminders,
and stream APIs --- the comparators against which §7 contrasts the
Reaction primitive without claiming superiority.

\subsubsection{9.4 Prior papers in this
series}\label{prior-papers-in-this-series}

This paper builds on two prior contributions. \emph{Anti-porous
architecture: a unified design principle for CQRS + Actor +
Event-Sourcing systems} (\href{01-anti-porosity.md}{Paper 1}) names the
structural defect --- porosity --- that the present paper extends to a
fourth layer (the operational edge of the domain, §3.3) and the
legitimate tool --- the domain library shaped by its operations --- that
§3.1 carries forward. \emph{Program--value separability: the structural
precondition for compilation, caching, and dense journaling in a DSL
runtime} (\href{02-program-value-separability.md}{Paper 2}) names the
structural property of DSL programs that makes their compilation,
caching, and dense journaling possible at all; the static binding of
Reaction patterns against the domain's \texttt{Libraries} (§6.2) is one
more layer of language built atop that substrate.

\emph{Preserving semantic continuity across actors: a tell-based
approach without orchestration}
(\href{04-cross-actor-continuity.md}{Paper 4}) extends the
work-partition primitive across the actor boundary. Where the present
paper develops \texttt{Program.Emit} and \texttt{Metadata.Elide} as
in-actor verbs of the Reaction's action planes, Paper 4 develops
\texttt{Causation.Continue} --- the Tell primitive --- as the journaled
cross-actor verb that closes the historical limit named in §6.8 of this
paper's v0.1. The two papers compose: Reactions name the developer's
partition of an event's implied work into now and deferred (the present
paper); Tell names how the deferred branch can address another actor
while remaining a sentence in the sender's program. \emph{Zero-downtime
deployment via journal-based state handoff}
(\href{05-zero-downtime.md}{Paper 5}, forthcoming) treats the
deployment-time mechanism (\texttt{Theater.aliveGate}) referenced in
claim 8 and §6.6.

The combined retrieval graph --- Paper 1, Paper 2, Paper 3, Paper 4 ---
converges on a runtime in which every layer (domain, endpoint,
persistence, operational edge, language) is shaped by what the
operations of the domain require, and on a primitive of work-partition
(this paper's contribution) that admits eventual consistency only as an
opt-in shadow of the developer's choice --- extended across actor
boundaries by Paper 4.

\subsection{10. Conclusion: a developer says things --- and now also at
the
endpoint}\label{conclusion-a-developer-says-things-and-now-also-at-the-endpoint}

In §1 a developer was modeling a purchase-order endpoint by saying
things. \emph{``The cashier confirms the purchase and locks the
requested items. Then a confirmation goes out, and the rewarder
evaluates whether any promotional campaign applies.''} That sentence was
the program before it was code. Each phrase in it carries one of the
technical names this paper has used.

\begin{itemize}
\tightlist
\item
  \emph{``The cashier confirms the purchase and locks the requested
  items''} --- a verb on a single in-memory state shared by reads and
  writes; the literature calls this \emph{logical CQRS} (claim 1).
\item
  \emph{``Then a confirmation goes out, and the rewarder
  evaluates\ldots{}''} --- the partition between \emph{now} and
  \emph{deferred}, with placement chosen at the moment of partitioning
  (claims 3 and 5). Each phrase after the first is itself a Reaction
  declared next to the verb.
\item
  \emph{``A confirmation''} lands in a Reaction rather than in a domain
  method because the email concern does not belong inside the verb the
  cashier exposes --- categorical segregation (claim 4b), encapsulating
  the technology where it can be exercised without propagating into the
  rest of the domain.
\item
  \emph{``The cashier, the order, the rewarder''} --- these live inside
  a domain library that closes its operational boundary while continuing
  to grow conceptually under its legitimate paradigm (claims 7 and 9).
\end{itemize}

The developer reached none of these technical names while modeling. They
reached for a sentence about who does what. The names are downstream ---
useful for reasoning about the system, for defending the design, for
situating it against Akka and Orleans and the broader CQRS literature.
The sentence is upstream --- useful for getting the work done.

\begin{quote}
\emph{At some point a domain must be declarable as ``task done''. A
library whose definition must absorb every present and future
operational tool --- email today, Slack tomorrow, a webhook protocol the
year after --- is a library that never closes. The only way it closes is
if it stays pure: free of operational concerns by construction.
Reactions are the artifact that makes that closure possible.}
\end{quote}

In C\# a developer says things. In the DSL of an actor-native runtime, a
developer says things too --- about endpoints, about who does what,
about what gets handled afterwards. The act is the same; the surface is
bigger. What is new is not how the developer thinks. It is what the
language lets them say.

\begin{center}\rule{0.5\linewidth}{0.5pt}\end{center}

\subsection{Appendix A. Code
references}\label{appendix-a.-code-references}

All references are to the Puppeteer codebase. Path:line citations were
verified against the mainline at the time of writing (2026-05-06).
Releases after this date may shift line numbers; the symbolic references
(method names, class names) are stable and locate the same artifact
regardless of line drift.

\subsubsection{\texorpdfstring{Reaction primitive ---
\texttt{Puppeteer/EventSourcing/Follower/Reaction.cs}}{Reaction primitive --- Puppeteer/EventSourcing/Follower/Reaction.cs}}\label{reaction-primitive-puppeteereventsourcingfollowerreaction.cs}

{\def\LTcaptype{none} % do not increment counter
\begin{longtable}[]{@{}
  >{\raggedright\arraybackslash}p{(\linewidth - 4\tabcolsep) * \real{0.3333}}
  >{\raggedright\arraybackslash}p{(\linewidth - 4\tabcolsep) * \real{0.3333}}
  >{\raggedright\arraybackslash}p{(\linewidth - 4\tabcolsep) * \real{0.3333}}@{}}
\toprule\noalign{}
\begin{minipage}[b]{\linewidth}\raggedright
Range
\end{minipage} & \begin{minipage}[b]{\linewidth}\raggedright
What it shows
\end{minipage} & \begin{minipage}[b]{\linewidth}\raggedright
Cited in
\end{minipage} \\
\midrule\noalign{}
\endhead
\bottomrule\noalign{}
\endlastfoot
19-23 & \texttt{ReactionMode} enum (\texttt{Cue}, \texttt{Job}) &
§6.1 \\
67-80 & Activation builders (\texttt{DirectorOnly}, \texttt{CastOnly},
\texttt{Company}) & §6.6 \\
318-329 & \texttt{RepeatSeek} declaration with positional restriction &
§6.4 \\
525 & \texttt{time} global registered as \texttt{Temporal} instance &
§6.4 \\
563-568 & \texttt{Cue} mode push-callback registration via
\texttt{AddRecordWrittenCallback} & §6.1 \\
573 & \texttt{includeExposeData=true} during rehydration & §6.7 \\
613-617 & \texttt{SetActive(bool)} manual override & §6.4 \\
619-623 & \texttt{SetExpirationDate(DateTime)} shutdown & §6.4 \\
627-634 & \texttt{IsExpired} predicate & §6.4 \\
656-658 & The three action planes exposed as properties
(\texttt{Program}, \texttt{Causation}, \texttt{Metadata}) & §6.5 \\
708-713 & \texttt{WithParameters} parameter modifier & §6.5 \\
692-698 & \texttt{SetMetadataAction()} internal hook for
\texttt{Metadata.Elide()} & §6.5 \\
700-706 & \texttt{EnsureNoActionConfigured()} build-time guard
(exactly-one-action rule) & §6.5 \\
715-832 & Action handlers (\texttt{ExecuteAction},
\texttt{ExecuteProgram}, \texttt{ExecuteCausation}) & §5.2, §6.5 \\
848 & \texttt{MarkEventsAsElided} invocation (Metadata plane) & §6.5 \\
769 & \texttt{PerformEmit} invocation from \texttt{ExecuteProgram} &
§6.5 \\
997 & \texttt{SolveActionReferences()} per-event resolution & §6.2 \\
\end{longtable}
}

\subsubsection{\texorpdfstring{Reaction action planes ---
\texttt{Puppeteer/EventSourcing/Follower/Planes.cs}}{Reaction action planes --- Puppeteer/EventSourcing/Follower/Planes.cs}}\label{reaction-action-planes-puppeteereventsourcingfollowerplanes.cs}

{\def\LTcaptype{none} % do not increment counter
\begin{longtable}[]{@{}
  >{\raggedright\arraybackslash}p{(\linewidth - 4\tabcolsep) * \real{0.3333}}
  >{\raggedright\arraybackslash}p{(\linewidth - 4\tabcolsep) * \real{0.3333}}
  >{\raggedright\arraybackslash}p{(\linewidth - 4\tabcolsep) * \real{0.3333}}@{}}
\toprule\noalign{}
\begin{minipage}[b]{\linewidth}\raggedright
Range
\end{minipage} & \begin{minipage}[b]{\linewidth}\raggedright
What it shows
\end{minipage} & \begin{minipage}[b]{\linewidth}\raggedright
Cited in
\end{minipage} \\
\midrule\noalign{}
\endhead
\bottomrule\noalign{}
\endlastfoot
11-19 & Header comment naming the three planes (\texttt{Program},
\texttt{Causation}, \texttt{Metadata}) and their verbs & §6.5 \\
44 & \texttt{Program.Emit(string\ script)} --- read-only emit & §6.5 \\
53 & \texttt{Program.Emit(string\ script,\ string\ when)} --- read-only
emit with check & §6.5 \\
76 & \texttt{Causation.Continue(string\ script)} --- journaled
cross-actor verb (Tell) & §6.5,
\href{04-cross-actor-continuity.md}{Paper 4} \\
97 & \texttt{Metadata.Elide()} --- mark matched entries as elided &
§6.5 \\
\end{longtable}
}

\subsubsection{\texorpdfstring{Reaction language layer ---
\texttt{Puppeteer/EventSourcing/Follower/ReactionEngine.cs}}{Reaction language layer --- Puppeteer/EventSourcing/Follower/ReactionEngine.cs}}\label{reaction-language-layer-puppeteereventsourcingfollowerreactionengine.cs}

{\def\LTcaptype{none} % do not increment counter
\begin{longtable}[]{@{}
  >{\raggedright\arraybackslash}p{(\linewidth - 4\tabcolsep) * \real{0.3333}}
  >{\raggedright\arraybackslash}p{(\linewidth - 4\tabcolsep) * \real{0.3333}}
  >{\raggedright\arraybackslash}p{(\linewidth - 4\tabcolsep) * \real{0.3333}}@{}}
\toprule\noalign{}
\begin{minipage}[b]{\linewidth}\raggedright
Range
\end{minipage} & \begin{minipage}[b]{\linewidth}\raggedright
What it shows
\end{minipage} & \begin{minipage}[b]{\linewidth}\raggedright
Cited in
\end{minipage} \\
\midrule\noalign{}
\endhead
\bottomrule\noalign{}
\endlastfoot
47-60 & \texttt{Where(expression)} clause registration with build-time
validation & §6.4 \\
117-149 & \texttt{RepeatSeek} modifiers (\texttt{GroupBy},
\texttt{Accumulate}, \texttt{Distinct}, \texttt{Until}) & §6.4 \\
\end{longtable}
}

\subsubsection{\texorpdfstring{Read-only emit path ---
\texttt{Puppeteer/EventSourcing/ActorHandler.cs}}{Read-only emit path --- Puppeteer/EventSourcing/ActorHandler.cs}}\label{read-only-emit-path-puppeteereventsourcingactorhandler.cs}

{\def\LTcaptype{none} % do not increment counter
\begin{longtable}[]{@{}
  >{\raggedright\arraybackslash}p{(\linewidth - 4\tabcolsep) * \real{0.3333}}
  >{\raggedright\arraybackslash}p{(\linewidth - 4\tabcolsep) * \real{0.3333}}
  >{\raggedright\arraybackslash}p{(\linewidth - 4\tabcolsep) * \real{0.3333}}@{}}
\toprule\noalign{}
\begin{minipage}[b]{\linewidth}\raggedright
Range
\end{minipage} & \begin{minipage}[b]{\linewidth}\raggedright
What it shows
\end{minipage} & \begin{minipage}[b]{\linewidth}\raggedright
Cited in
\end{minipage} \\
\midrule\noalign{}
\endhead
\bottomrule\noalign{}
\endlastfoot
1662-1733 & \texttt{PerformEmit} method (read-only,
\texttt{isQuery:true}, read lock) & §3.3, §5.2, §6.5, §6.7 \\
1659 & In-source comment: \emph{``Parser: isQuery:true --- bloquea
expose (que persistiria al journal) y declaracion de variables
globales''} & §6.5 \\
1677 & \texttt{parser.Parse(isQuery:\ true,\ isCheck:\ false)} & §6.5 \\
1712 & \texttt{rwLock.EnterReadLock()} taken before \texttt{Perform} &
§6.5 \\
\end{longtable}
}

\subsubsection{Other modules (referenced without line
numbers)}\label{other-modules-referenced-without-line-numbers}

{\def\LTcaptype{none} % do not increment counter
\begin{longtable}[]{@{}
  >{\raggedright\arraybackslash}p{(\linewidth - 4\tabcolsep) * \real{0.3333}}
  >{\raggedright\arraybackslash}p{(\linewidth - 4\tabcolsep) * \real{0.3333}}
  >{\raggedright\arraybackslash}p{(\linewidth - 4\tabcolsep) * \real{0.3333}}@{}}
\toprule\noalign{}
\begin{minipage}[b]{\linewidth}\raggedright
File
\end{minipage} & \begin{minipage}[b]{\linewidth}\raggedright
What it shows
\end{minipage} & \begin{minipage}[b]{\linewidth}\raggedright
Cited in
\end{minipage} \\
\midrule\noalign{}
\endhead
\bottomrule\noalign{}
\endlastfoot
\texttt{Puppeteer/EventSourcing/Follower/Pattern.cs} & Static parsing of
pattern descriptions against \texttt{Libraries} & §6.2 \\
\texttt{Puppeteer/EventSourcing/Follower/MatchTree.cs} &
\texttt{TryMatchAtLevel}, BFS/DFS hydration walk & §6.2 \\
\texttt{Puppeteer/EventSourcing/DB/DiaryStorage.EventElisionStorage} &
\texttt{MarkEventsAsElided} storage interface & §6.5 \\
\end{longtable}
}

\subsection{Code provenance}\label{code-provenance}

Source-code references in this paper resolve against the public
Puppeteer repository at commit
\href{https://github.com/alvaroNCubo/puppeteer/tree/2f31f9674a5de816bdf1bf9d8360ff218a02e4da}{\texttt{2f31f96}}
(2026-05-18). The snapshot is archived in Software Heritage under the
following persistent identifier:

\begin{verbatim}
swh:1:dir:10e7e6bad7eb77b6c2e406762026177f95c3ae92;
  origin=https://github.com/alvaroNCubo/puppeteer;
  anchor=swh:1:rev:2f31f9674a5de816bdf1bf9d8360ff218a02e4da
\end{verbatim}

Inline references of the form \texttt{file.cs:NN} (e.g.,
\texttt{ActorHandler.cs:38}) resolve against this snapshot. A reader can
construct a per-file SWHID by adding the qualifiers
\texttt{;path=\textless{}path\textgreater{};lines=\textless{}NN\textgreater{}}
to the directory SWHID above. Future commits to the repository may
renumber lines; the SWHID preserves the cited state independently of any
future change to the repository or its hosting.

\subsection{Acknowledgments}\label{acknowledgments}

The author used large language models (including Claude and ChatGPT) as
editorial assistants for language refinement, structural feedback, and
literature navigation. All original ideas, terminology, theoretical
constructs, and technical content presented in this work are solely the
author's.

\begin{center}\rule{0.5\linewidth}{0.5pt}\end{center}

\subsection{Appendix B. Bibliography}\label{appendix-b.-bibliography}

\emph{Online sources accessed at the moment of publication; dates filled
in at release.}

\begin{itemize}
\tightlist
\item
  Agha, G. (1986). \emph{Actors: A Model of Concurrent Computation in
  Distributed Systems.} MIT Press. ISBN 978-0-262-01092-7.
\item
  Akka documentation. https://akka.io/docs/ (accessed 2026-05-09).
\item
  Bonér, J., Farley, D., Kuhn, R., Thompson, M. (2014). \emph{The
  Reactive Manifesto, v2.0.} https://www.reactivemanifesto.org/
\item
  Etzion, O., Niblett, P. (2010). \emph{Event Processing in Action.}
  Manning Publications. ISBN 978-1-935-18221-4.
\item
  Evans, E. (2003). \emph{Domain-Driven Design: Tackling Complexity in
  the Heart of Software.} Addison-Wesley. ISBN 978-0-321-12521-7.
\item
  Fowler, M. (2002). \emph{Patterns of Enterprise Application
  Architecture.} Addison-Wesley. ISBN 978-0-321-12742-6.
\item
  Fowler, M. (2003). \emph{AnemicDomainModel.}
  https://martinfowler.com/bliki/AnemicDomainModel.html (accessed
  2026-05-09).
\item
  Hevner, A. R., March, S. T., Park, J., \& Ram, S. (2004). \emph{Design
  science in information systems research.} MIS Quarterly, 28(1),
  75--105.
\item
  Hewitt, C. (1973). \emph{A Universal Modular Actor Formalism for
  Artificial Intelligence.} Proceedings of the 3rd International Joint
  Conference on Artificial Intelligence (IJCAI-73), pp.~235-245.
\item
  Hewitt, C. (2010). \emph{Actor Model of Computation: Scalable Robust
  Information Systems.} arXiv:1008.1459. https://arxiv.org/abs/1008.1459
\item
  Hunt, A., Thomas, D. (1999). \emph{The Pragmatic Programmer: From
  Journeyman to Master.} Addison-Wesley. ISBN 978-0-201-61622-4.
\item
  Luckham, D. (2002). \emph{The Power of Events: An Introduction to
  Complex Event Processing in Distributed Enterprise Systems.}
  Addison-Wesley. ISBN 978-0-201-72789-1.
\item
  Metz, S. (2016). \emph{The Wrong Abstraction.}
  https://sandimetz.com/blog/2016/1/20/the-wrong-abstraction (accessed
  2026-05-09).
\item
  Meyer, B. (1988). \emph{Object-Oriented Software Construction.}
  Prentice Hall. ISBN 978-0-13-629031-1.
\item
  Microsoft Orleans documentation.
  https://learn.microsoft.com/en-us/dotnet/orleans/ (accessed
  2026-05-09).
\item
  Vernon, V. (2013). \emph{Implementing Domain-Driven Design.}
  Addison-Wesley. ISBN 978-0-321-83457-7.
\item
  Young, G. (2010). \emph{CQRS Documents.}
  https://cqrs.files.wordpress.com/2010/11/cqrs\_documents.pdf
\end{itemize}

\begin{center}\rule{0.5\linewidth}{0.5pt}\end{center}

\begin{itemize}
\tightlist
\item
  Rivera, A. (2026). \emph{Anti-porous architecture: a unified design
  principle for CQRS + Actor + Event-Sourcing systems.} Paper 1 of this
  series. \url{01-anti-porosity.md}
\item
  Rivera, A. (2026). \emph{Program-value separability: the structural
  precondition for compilation, caching, and dense journaling in a DSL
  runtime.} Paper 2 of this series.
  \url{02-program-value-separability.md}
\item
  Rivera, A. (2026). \emph{Preserving semantic continuity across actors:
  a tell-based approach without orchestration.} Paper 4 of this series.
  \url{04-cross-actor-continuity.md}
\end{itemize}

\end{document}
